\documentclass[11pt,a4paper]{article}
\usepackage[margin=22mm,headheight=18pt]{geometry}
\usepackage{amsmath,amssymb}
\usepackage{graphicx}
\usepackage{xcolor}
\usepackage{booktabs,tabularx,array}
\usepackage{listings}
\usepackage{fancyhdr}
\usepackage{hyperref}
\definecolor{ink}{HTML}{193047}
\definecolor{teal}{HTML}{008A8A}
\definecolor{pale}{HTML}{F0F5F7}
\hypersetup{colorlinks=true,linkcolor=teal,urlcolor=teal,pdftitle={راهنمای آموزشی پروژه CommitJL},pdfauthor={Project research workspace}}
\usepackage{xepersian}
\settextfont[Path=C:/Windows/Fonts/,BoldFont=tahomabd.ttf]{tahoma.ttf}
\setlatintextfont{Latin Modern Roman}
\setmonofont[Path=C:/Windows/Fonts/,BoldFont=consolab.ttf,ItalicFont=consolai.ttf,BoldItalicFont=consolaz.ttf]{consola.ttf}
\linespread{1.22}
\setlength{\parindent}{0pt}
\setlength{\parskip}{7pt}
\setlength{\emergencystretch}{3em}
\pagestyle{fancy}
\fancyhf{}
\fancyhead[R]{\small راهنمای آموزشی پروژه}
\fancyhead[L]{\small\lr{CommitJL}}
\fancyfoot[C]{\thepage}
\renewcommand{\headrulewidth}{0.3pt}
\lstset{language=Python,basicstyle=\ttfamily\footnotesize,breaklines=true,columns=fullflexible,
  keepspaces=true,showstringspaces=false,backgroundcolor=\color{pale},frame=single,
  rulecolor=\color{pale},tabsize=4,aboveskip=8pt,belowskip=8pt}
\newcommand{\lesson}[1]{\clearpage\section{#1}}
\newcommand{\codefile}[1]{\begin{latin}\lstinputlisting{#1}\end{latin}}
\newcommand{\chart}[3][0.96]{\begin{center}\includegraphics[width=#1\linewidth]{#2}\end{center}{\small #3}\par}
\newcommand{\takeaway}[1]{\par\textbf{چیزی که باید یادت بماند:} #1\par}
\newcommand{\status}[1]{\textbf{وضعیت این بخش:} #1\par}
\input{proof_macros.tex}
\begin{document}
\begin{titlepage}
\vspace*{18mm}
{\Huge\bfseries از یک جدول داده تا یک اثبات بدون افشای داده}\par
\vspace{12mm}
{\LARGE راهنمای قدم‌به‌قدم مقاله و پروژه}\par
\vspace{5mm}
{\LARGE\lr{CommitJL / Zero-Knowledge Clustering}}\par
\vspace{14mm}
{\large برای کسی که از خوشه‌بندی، کاهش بُعد و دانش صفر چیزی نمی‌داند}\par
\vspace{12mm}
این راهنما توضیح می‌دهد مسئله چیست، روش ما چه کاری انجام می‌دهد، چرا PCA رقیب مهمی است، داده‌های واقعی چه شکلی دارند و نتایجی که تا امروز گرفته‌ایم واقعاً چه می‌گویند. مثال‌های عددی، تصویر، کد و یک اثبات کوچک واقعی در متن آمده‌اند.
\vspace{8mm}
\textbf{تاریخ بررسی:} ۴ اکتبر ۲۰۲۶\par
\textbf{مبنای متن:} پیش‌نویس مقاله، کد فعلی مخزن، نتایج ران تشخیصی اخیر و مثال‌های آموزشی اجراشده.\par
\textbf{این فایل یک راهنمای آموزشی است، نه مقاله نهایی و نه ادعای تکمیل همه بنچمارک‌ها.}\par
\vfill
\textcolor{teal}{\large ایده نظری، قابلیت پیاده‌سازی‌شده و نتیجه اندازه‌گیری‌شده سه چیز متفاوت‌اند. در این متن آن‌ها را جدا نگه می‌داریم.}
\end{titlepage}
\begingroup\small\linespread{1.08}\selectfont\tableofcontents\endgroup

\lesson{داستان کار از کجا شروع می‌شود؟}
فرض کن یک بیمارستان برای هر بیمار یک جدول بزرگ دارد: هزاران عدد از آزمایش‌ها، تصاویر یا اندازه‌گیری‌های زیستی. یک پژوهشگر می‌خواهد بفهمد بیماران از نظر این ویژگی‌ها به چند گروه شبیه تقسیم می‌شوند. بیمارستان نمی‌خواهد اطلاعات تک‌تک افراد را منتشر کند؛ پژوهشگر هم نمی‌خواهد صرفاً به جمله «ما محاسبه کردیم و درست است» اعتماد کند.

پس دو خواسته داریم: \textbf{داده خصوصی بماند} و \textbf{درستی یک محاسبه مشخص قابل بررسی باشد}. ابزار دانش صفر می‌تواند به خواسته دوم بدون افشای شاهد خصوصی کمک کند. اما وقتی محاسبه روی یک جدول خیلی بزرگ انجام شود، ساخت اثبات می‌تواند گران و پرمصرف شود.

ایده ما این است: قبل از انجام مقایسه‌های فاصله، هر ردیف را از یک بردار بلند به یک بردار کوتاه تبدیل کنیم. مثلاً به‌جای ۷۸۴ عدد برای تصویر، ۶۴ عدد داشته باشیم. اگر این کوتاه‌کردن هندسه را به اندازه کافی حفظ کند، انتساب به نزدیک‌ترین مرکز خوشه هم ممکن است حفظ شود. سپس درستی محاسبه کوتاه‌شده را اثبات کنیم.

این کار مثل خلاصه‌کردن یک پرونده با چند شاخص است: خلاصه سریع‌تر خوانده می‌شود، ولی اگر چیز مهمی حذف شود تصمیم عوض می‌شود. تمام مسئله این است که بفهمیم \textbf{چه خلاصه‌ای، روی چه داده‌ای، با چه میزان خطا و با چه هزینه اثباتی قابل دفاع است}.

\textbf{یک ظرافت مهم:} فشرده‌کردن داده، خودش ابزار حفظ حریم خصوصی نیست. ممکن است از داده فشرده‌شده هم اطلاعات استخراج شود. در پروتکل موردنظر، داده خصوصی و تصویرش داخل شاهد اثبات قرار می‌گیرند؛ نقش کاهش بُعد، کم‌کردن هزینه محاسبه و نقش ZK، پنهان‌کردن شاهد در اثبات است.

\takeaway{ما نمی‌خواهیم صرفاً خوشه‌بندی سریع‌تر بسازیم. می‌خواهیم یک محاسبه خوشه‌ای را با داده خصوصی، قابل راستی‌آزمایی و تا حد ممکن کم‌هزینه کنیم.}

\lesson{یک جدول داده دقیقاً چیست؟}
یک دیتاست را به شکل ماتریس \(X\) می‌نویسیم. هر سطر یک نمونه و هر ستون یک ویژگی است. برای سه مشتری با دو ویژگی می‌توانیم چنین جدولی داشته باشیم:
\[
 X=\begin{bmatrix}20&2\\22&3\\60&12\end{bmatrix}.
\]
ستون اول می‌تواند سن و ستون دوم تعداد خرید باشد. این فقط یک مثال است؛ در MNIST هر ستون یک پیکسل است و در داده متنی هر ستون وزن یک واژه.

\begin{latin}
\begin{center}\begin{tabular}{cll}
\toprule Symbol & Meaning & Example \\\midrule
\(n\) & Number of rows & 1,000 images \\
\(d\) & Original features & 784 pixels \\
\(m\) & Projected features & 64 coordinates \\
\(k\) & Number of centers & 5 clusters \\
\(q\) & Nonzeros per sparse column & 2 or 8 \\
\(S\) & Fixed-point scale & \(2^{16}\) \\
\(B\) & Integer coordinate bound & 4,096 \\
\(Q\) & Allowed retry budget & 1 or 8 \\
\bottomrule\end{tabular}\end{center}
\end{latin}

\textbf{نباید \(q\) را با تعداد بیت دقت اشتباه بگیریم.} در ریاضیات روش و کد اصلی، \(q\) تعداد درایه‌های ناصفر در هر ستون ماتریس تنک است. آزمایش دقت از برچسب \lr{precision\_bits} و مقیاس \(S\) استفاده می‌کند. همچنین \(Q\) بزرگ تعداد تلاش‌های مجاز در آزمایش grinding است. این سه مفهوم کاملاً جدا هستند؛ برخی درخواست‌ها و نام‌گذاری‌های قبلی آن‌ها را مخلوط کرده‌اند.

وقتی از \(d=784\) به \(m=64\) می‌رویم، تعداد مختصات حدود ۹۱٫۸ درصد کم می‌شود. این عدد به‌تنهایی به معنی ۹۱٫۸ درصد کاهش کل زمان، حافظه یا قید اثبات نیست؛ بخش‌هایی از پروتکل هنوز به \(nd\) عدد اصلی دسترسی دارند.

\takeaway{بُعد، تعداد ویژگی هر نمونه است؛ تعداد نمونه، تعداد ردیف‌هاست. کم‌کردن بُعد با کم‌کردن تعداد تصاویر فرق دارد.}

\lesson{خوشه‌بندی: گروه‌بندی شباهت، نه تشخیص جواب درست}
در طبقه‌بندی معمولاً جواب آموزشی داریم: این تصویر «۸» است، آن تصویر «۵». مدل از نمونه‌های برچسب‌دار یاد می‌گیرد. در خوشه‌بندی می‌گوییم نمونه‌های شبیه را کنار هم بگذار، بدون این‌که هنگام تصمیم‌گیری برچسب واقعی را به الگوریتم بدهیم.

در \lr{KMeans}، \(k\) مرکز تعریف می‌کنیم. هر نقطه به نزدیک‌ترین مرکز می‌رود؛ بعد هر مرکز به میانگین اعضای خوشه‌اش منتقل می‌شود. این دو مرحله چند بار تکرار می‌شوند. هدف معمول، کم‌کردن مجموع فاصله‌های مربعی است:
\[
 D(x,c)=\sum_{t=1}^{d}(x_t-c_t)^2,
 \qquad \Phi(X,C)=\sum_{i=1}^{n}\min_j D(x_i,c_j).
\]
برای \(x=(2,2)\)، \(c_1=(1,1)\) و \(c_2=(5,5)\)، فاصله‌های مربعی ۲ و ۱۸ هستند، پس مرکز اول انتخاب می‌شود. نیازی به گرفتن ریشه مربع نیست؛ ترتیب فاصله‌ها عوض نمی‌شود.

اگر اعضای یک خوشه \((1,1)\) و \((3,3)\) باشند، مرکز جدید \((2,2)\) است. این «به‌روزرسانی مرکز» با «انتساب به مرکز ثابت» فرق دارد.

\textbf{در کد فعلی ما:} ابتدا KMeans در فضای اصلی اجرا می‌شود و مرکزها به‌دست می‌آیند. سپس همان نقطه‌ها و همان مرکزها به فضای کوتاه برده می‌شوند و انتساب دوباره محاسبه می‌شود. مدار فعلی درستی انتساب به مرکزهای داده‌شده را بررسی می‌کند؛ همه مراحل آموزش KMeans را اثبات نمی‌کند.

عدد \(k=5\) یعنی پنج گروه شباهت، نه لزوماً پنج مفهوم انسانی مشخص. شماره خوشه‌ها هم قراردادی است: عوض‌کردن نام خوشه صفر و یک، کیفیت خوشه‌بندی را تغییر نمی‌دهد.

\takeaway{اثبات نزدیک‌ترین مرکز بودن، با اثبات این‌که مرکزها بهترین مرکزهای ممکن‌اند یا مدل معنای واقعی کلاس‌ها را یاد گرفته، یکسان نیست.}

\lesson{MNIST و سؤال زوج و فرد}
MNIST مجموعه‌ای از تصویرهای خاکستری ارقام دست‌نویس صفر تا نه است. هر تصویر \(28\times28\) پیکسل دارد. مقدار هر پیکسل در نسخه بارگذاری‌شده یک عدد از صفر تا ۲۵۵ است. تصویر تخت می‌شود و یک بردار با ۷۸۴ مؤلفه می‌سازد:
\[
\text{image}_{28\times28}\longrightarrow x\in\mathbb{R}^{784}.
\]
برچسب «۸» بیرون این بردار نگه داشته می‌شود. خود تصویر یک عدد ۸ نیست؛ ۷۸۴ شدت روشنایی است که شکل رقم را می‌سازند.

\chart{mnist_samples.png}{شش تصویر واقعی از زیرمجموعه بارگذاری‌شده. عنوان هر تصویر برچسب اصلی آن است، نه خروجی خوشه‌بندی.}

می‌توانیم \(k=2\) بگذاریم، ولی KMeans به‌طور خودکار «زوج و فرد» را نمی‌فهمد. ممکن است تصاویر حلقه‌دار را از تصاویر کشیده جدا کند، یا اندازه و ضخامت خط روی تقسیم اثر بگذارد. رقم‌های ۳ و ۸ شاید از نظر پیکسل شبیه‌تر از ۲ و ۸ باشند، با این‌که اولی فرد و دومی زوج است.

برای آزمون مشخص زوج/فرد، می‌توان برچسب ارزیابی را \(y\bmod2\) تعریف کرد و بعد نتیجه خوشه‌بندی را با آن مقایسه کرد. ولی این فقط ارزیابی یک هدف انسانی است؛ تغییر \(k\) به ۲ تضمین نمی‌کند چنین ساختاری قابل کشف باشد. اگر واقعاً هدف تشخیص زوج و فرد باشد، یک روش طبقه‌بندی آموزش‌دیده معمولاً سؤال مناسب‌تری را پاسخ می‌دهد.

در ران فعلی غالباً \(k=5\) است، در حالی که MNIST ده کلاس رقم دارد. بنابراین معیار اصلی ما «چقدر انتساب بعد از کاهش بُعد به انتساب قبل از آن شبیه است» است، نه دقت تشخیص ده رقم.

\takeaway{برچسب واقعی برای قضاوت پس از محاسبه مفید است؛ الگوریتم خوشه‌بندی در این آزمایش از آن برای آموزش استفاده نمی‌کند.}

\lesson{کاهش بُعد: چه چیزی را حذف می‌کنیم؟}
فرض کن برای هر خانه متراژ را هم با مترمربع و هم با فوت‌مربع ثبت کرده‌ای. دو ستون تقریباً یک اطلاعات را می‌گویند؛ نگه‌داشتن هر دو ممکن است لازم نباشد. در یک تصویر هم بسیاری از پیکسل‌ها هم‌بسته‌اند: ضخیم‌ترشدن خط، چند پیکسل کنار هم را تغییر می‌دهد.

کاهش بُعد تلاش می‌کند تعداد کمتر مختصات را جایگزین تعداد بیشتر کند. در کاهش خطی یک ماتریس \(P\) داریم:
\[
 z=Px,\qquad P\in\mathbb{R}^{m\times d},\qquad m<d.
\]
هر مختصه جدید یک ترکیب از مختصات قدیمی است. مثلاً \(z_1=x_1+x_2-x_4\). بنابراین ۶۴ مؤلفه PCA معمولاً ۶۴ پیکسل انتخاب‌شده نیستند؛ ترکیب‌هایی از همه پیکسل‌ها هستند.

وقتی \(m<d\) باشد، در حالت کلی اطلاعاتی از دست می‌رود. نمی‌توان انتظار داشت همه بردارهای ممکن و همه فاصله‌های آن‌ها دقیقاً حفظ شوند. هدف، حفظ ساختار موردنیاز روی یک مجموعه مشخص از داده‌ها با خطای قابل کنترل است.

دو سؤال مستقل می‌پرسیم: آیا شکل هندسی و تصمیم خوشه‌ای حفظ می‌شود؟ و آیا هزینه اجرای خصوصی و اثبات پایین می‌آید؟ یک تبدیل ممکن است در سؤال اول عالی باشد ولی در سؤال دوم گران یا سخت‌احراز باشد.

اگر \(m=d\) باشد، حتی یک ماتریس تصادفی هم «کاهش بُعد» محسوب نمی‌شود. اگر تعداد سطرهای دیتاست را از ۷۰هزار به هزار کاهش داده باشیم، این «زیرنمونه‌گیری دیتاست» است، نه همان کاهش بُعد هر نمونه. در گزارش باید هر دو را شفاف بنویسیم.

\takeaway{ما هم هزینه اطلاعاتی کاهش بُعد را داریم و هم هزینه اثبات آن را. بهتر بودن در یکی، برتری در دیگری را تضمین نمی‌کند.}

\lesson{PCA از صفر: بهترین جهت برای دیدن پراکندگی}
PCA مخفف تحلیل مؤلفه‌های اصلی است. تصور کن از ابر نقطه‌ها عکس می‌گیری: از بعضی زاویه‌ها نقطه‌ها روی هم می‌افتند؛ از بعضی زاویه‌ها پراکندگی بیشتری دیده می‌شود. PCA جهت‌هایی پیدا می‌کند که پراکندگی داده را بیشترین مقدار نگه می‌دارند. جهت‌ها عمود بر هم‌اند.

ابتدا میانگین هر ستون را کم می‌کند، سپس با تجزیه ماتریسی مانند SVD جهت‌ها را یاد می‌گیرد. برای یک مثال کاملاً ساده:
\[
 X=\begin{bmatrix}1&1\\2&2\\3&3\end{bmatrix},\quad
 \mu=(2,2),\quad
 X-\mu=\begin{bmatrix}-1&-1\\0&0\\1&1\end{bmatrix}.
\]
همه نقطه‌ها روی خط \(x_2=x_1\) هستند. جهت اصلی واحد برابر \(w=(1,1)/\sqrt2\) است و هر نقطه یک نمره جدید می‌گیرد:
\[
 z=(x-\mu)\cdot w,
 \qquad Z=\begin{bmatrix}-\sqrt2\\0\\\sqrt2\end{bmatrix}.
\]
این‌جا با یک عدد به‌جای دو عدد می‌توان نقطه‌ها را دقیق بازسازی کرد، چون تمام داده روی یک خط قرار دارد. در داده واقعی معمولاً چنین بازسازی بی‌نقصی نداریم.

\chart{toy_pca.png}{مثال اجراشده: یک جهت برای ثبت جای سه نقطه روی خط کافی است.}

ماتریس کوواریانس نمونه این مثال \(\begin{bmatrix}1&1\\1&1\end{bmatrix}\) است. جهت اصلی مقدار ویژه ۲ دارد و جهت عمود مقدار ویژه صفر؛ یعنی در جهت دوم هیچ پراکندگی‌ای وجود ندارد. توضیح شهودی همین است: ستون دوم چیز مستقلی به ستون اول اضافه نمی‌کند.

\takeaway{PCA از خود داده یاد می‌گیرد کدام جهت‌ها مهم‌ترند؛ SparseJL چنین آموزشی ندارد.}

\lesson{PCA در کد: fit، transform و بازسازی}
این مثال واقعاً با scikit-learn اجرا شده است:
\begin{latin}\begin{lstlisting}
import numpy as np
from sklearn.decomposition import PCA

X = np.array([[1., 1.], [2., 2.], [3., 3.]])
pca = PCA(n_components=1, svd_solver="full")
Z = pca.fit_transform(X)
X_reconstructed = pca.inverse_transform(Z)
assert np.allclose(X_reconstructed, X)
print(pca.mean_, pca.components_, Z)
\end{lstlisting}\end{latin}
\codefile{toy_pca_output.txt}
\lr{fit} یعنی جهت‌ها و میانگین را از داده یاد بگیر؛ \lr{transform} یعنی با جهت‌های یادگرفته‌شده نمونه را کوتاه کن؛ \lr{inverse\_transform} یعنی بازسازی تقریبی را در فضای اصلی بساز. انجام PCA روی تنها یک تصویر، یک مدل عمومی از تغییرات بین تصویرها به ما نمی‌دهد؛ برای یادگیری جهت‌ها به مجموعه تصویرها نیاز داریم.

علامت جهت PCA می‌تواند برعکس شود. اگر نرم‌افزار \((-1,-1)/\sqrt2\) بدهد، نمره‌ها هم برعکس می‌شوند و فاصله‌ها همان می‌مانند. این خطا نیست. همچنین مرکزکردن با استانداردکردن ستون‌ها فرق دارد: PCA معمولی میانگین را کم می‌کند، اما خودکار هر ستون را به واریانس یک تبدیل نمی‌کند [۱].

\textbf{چه چیزی برای ZK مهم است؟} اگر ماتریس جهت‌ها از قبل عمومی و ثابت باشد، اعمال آن یک محاسبه خطی مشخص است. اگر بخواهیم اثبات کنیم خود جهت‌ها دقیقاً PCA داده خصوصی هستند، باید یادگیری و ویژگی‌های تجزیه هم احراز شود؛ مسئله دیگری است. مقایسه منصفانه نباید برای روش ما فقط اعمال ماتریس و برای PCA کل آموزش را حساب کند، مگر این‌که این دو تعریف متفاوت را آشکارا بیان کنیم.

\takeaway{«PCA را روی داده اعمال کردیم» و «درستی یادگیری PCA را هم اثبات کردیم» دو ادعای متفاوت‌اند.}

\lesson{PCA روی داده واقعی MNIST}
در این راهنما از همان بارگذار محلی، هزار تصویر واقعی با seed صفر گرفته‌ایم. یک وصله کوچک \(4\times6\) از اولین تصویر را هم استخراج کرده‌ایم؛ این اعداد واقعاً پیکسل‌اند، نه مثال ساختگی. سپس PCA با ۶۴ مؤلفه اجرا شده است.
\begin{latin}\begin{lstlisting}
from data.loaders import load_dataset
from sklearn.decomposition import PCA

data = load_dataset("mnist", n_max=1000, seed=0)
X = data["X"]  # (1000, 784), raw pixels in [0, 255]
model = PCA(n_components=64, svd_solver="full")
Z = model.fit_transform(X)  # (1000, 64)
print(model.explained_variance_ratio_.sum())
patch = X[0].reshape(28, 28)[10:14, 10:16]
\end{lstlisting}\end{latin}
\codefile{mnist_pca_output.txt}
۶۴ مؤلفه حدود \textbf{۸۷٫۴۶ درصد واریانس} این زیرمجموعه را حفظ کردند. این عدد «۸۷٫۴۶ درصد دقت تشخیص رقم» نیست. واریانس، میزان پراکندگی مختصات است؛ نگه‌داشتن آن لزوماً همه جزئیات مهم برای یک تصمیم را حفظ نمی‌کند.

\chart{mnist_reconstruction.png}{تصویر اول و بازسازی آن با ۸، ۱۶ و ۶۴ مؤلفه از مدل یکسان. محوشدن یا حذف جزئیات، اثر کاهش بُعد است. این تصویر نتیجه آموزشی است، نه یک ردیف جدید از بنچمارک مقاله.}

هزار تصویر یک زیرمجموعه از مخزن ۷۰هزار تصویر است. این مثال و ران تشخیصی، ارزیابی آموزش/آزمون جداگانه برای طبقه‌بندی نیستند؛ درباره حفظ ساختار روی همین مجموعه صحبت می‌کنند.

\lesson{بعد از PCA یا SparseJL چه محاسبه‌ای انجام می‌دهیم؟}
ابتدا روی فضای اصلی KMeans با پنج مرکز و seed صفر اجرا می‌کنیم. سپس هم نمونه‌ها و هم مرکزها را با یک تبدیل به فضای کوتاه می‌بریم. تصمیم جدید، نزدیک‌ترین مرکز \textbf{پروجکت‌شده} است؛ KMeans جدیدی را در فضای کوتاه از نو آموزش نمی‌دهیم.
\begin{latin}\begin{lstlisting}
from sklearn.cluster import KMeans
from commitjl.clustering import nearest_centroid_assign
from experiments.pipeline import build_projection

km = KMeans(n_clusters=5, n_init=1, random_state=0).fit(X)
C = km.cluster_centers_
full_labels = nearest_centroid_assign(X, C)
pca_labels = nearest_centroid_assign(
    model.transform(X), model.transform(C))
P = build_projection("sparse_jl", X, 64, 2, 0, 65536)["real"]
sparse_labels = nearest_centroid_assign(X @ P.T, C @ P.T)
\end{lstlisting}\end{latin}
این قطعه ادامه مثال MNIST صفحه قبل است و از همان \(X\) و مدل PCA استفاده می‌کند. تابع ساخت projection نمایش حقیقی ماتریس را برمی‌گرداند. خروجی واقعی مقایسه در ادامه است:
\codefile{assignments_output.txt}
در این \textbf{یک مثال با یک seed}، ARI بین PCA و مرجع اصلی ۱ و بین SparseJL و مرجع حدود ۰٫۷۰۰ بود. این را نباید به جای میانگین ده seed در نتایج RQ1 بگذاریم. رقم واقعی ۸ هم ممکن است در خوشه شماره ۱ باشد؛ شماره خوشه نام رقم نیست.

در برخی مسیرهای کد، PCA به صورت \(XW^T\) به‌جای \((X-\mu)W^T\) اعمال می‌شود. اگر همان تبدیل برای مرکزها هم استفاده شود، یک انتقال مشترک \(\mu W^T\) حذف می‌شود و فاصله نقطه تا مرکز ثابت می‌ماند. این هم‌ارزی در مثال بررسی شده؛ برای بازسازی تصویر یا تفسیر نمره PCA همچنان میانگین مهم است.

\takeaway{مقایسه فعلی اثر تغییر هندسه بر تصمیم‌های مرکزهای ثابت را جدا می‌کند؛ نتیجه آموزش کامل در فضای کاهش‌یافته نیست.}

\lesson{baselineها: رقیب‌ها هر کدام چه می‌کنند؟}
Baseline یک روش مرجع است که کمک می‌کند بفهمیم روش جدید چه چیزی اضافه کرده است. بدون آن، حتی یک عدد خوب هم معلوم نمی‌کند پیشرفت داشته‌ایم یا مسئله ذاتاً آسان بوده است.

\textbf{فضای اصلی یا Full:} هیچ کاهش بُعدی انجام نمی‌شود. مرجع انتساب‌ها و هزینه هندسی است. طبیعی است ARI آن نسبت به خودش ۱ باشد؛ این به معنی تشخیص بی‌نقص کلاس واقعی نیست.

\textbf{PCA:} جهت‌های مهم را از داده یاد می‌گیرد. معمولاً برای داده‌های دارای هم‌بستگی و ساختار کم‌بُعد، حفظ ساختار خوبی دارد. ماتریس آن غالباً متراکم است. هزینه fit و هزینه transform باید جدا گزارش شوند.

\textbf{Gaussian:} هر درایه ماتریس تصادفی از توزیع نرمال با میانگین صفر و واریانس \(1/m\) است. همه مختصات معمولاً در همه خروجی‌ها دخالت دارند. احتمال حفظ فاصله‌ها، موضوع نظریه JL است؛ یک draw خاص تضمین بی‌خطا ندارد.

\textbf{Rademacher:} درایه‌ها با احتمال برابر \(+1/\sqrt m\) یا \(-1/\sqrt m\) هستند. انتخاب علامت‌ها ساده است، ولی ماتریس همچنان متراکم است.

\textbf{Achlioptas:} در نسخه موجود، احتمال صفر دو سوم و احتمال هر علامت یک ششم است؛ ضریب حقیقی مقیاس \(\sqrt{3/m}\) دارد. تنک است، اما تعداد ناصفرهای هر ستون تصادفی است.

\textbf{SparseJL ما:} هر ستون دقیقاً \(q\) ناصفر با علامت‌های مثبت یا منفی دارد. تعداد کل ناصفرها دقیقاً \(dq\) است. این ساختار برای نمایش صحیح و احراز تنکی مناسب است.

\textbf{Coordinate sampling:} چند ستون انتخاب می‌شوند. بسیار ساده است، ولی شاید ستون‌های مهم را حذف کند. در پیاده‌سازی موجود مقیاس جبرانی \(\sqrt{d/m}\) ندارد؛ بنابراین baseline ساده فعلی از نظر نرمال‌سازی فاصله‌ها با JL استاندارد یکسان نیست.

\textbf{Pisot:} baseline قطعی/شبه‌تصادفی الهام‌گرفته از تبدیل‌های بتا. بازسازی فعلی ما با artifact مقاله مرجع تطبیق کامل و مستقل نشده است؛ آن را بازتولید تأییدشده نتایج مقاله دیگر معرفی نمی‌کنیم.

\takeaway{یک روش می‌تواند برای دقت بومی بهترین و برای یک پروتکل اثبات خاص نامناسب باشد. برتری نهایی باید با تعریف یکسان مسئله اندازه‌گیری شود.}

\lesson{ماتریس‌های تصادفی و یک نمونه SparseJL با دست}
این کد شکل رایج baselineهای متراکم و Achlioptas را نشان می‌دهد. مثال آموزشی است و جایگزین فایل تولیدکننده رسمی پروتکل نیست:
\begin{latin}\begin{lstlisting}
rng = np.random.default_rng(7)
m, d = 2, 4
gaussian = rng.normal(0, 1 / np.sqrt(m), (m, d))
rademacher = rng.choice([-1, 1], (m, d)) / np.sqrt(m)
achlioptas = rng.choice([-1, 0, 1], (m, d),
                        p=[1/6, 2/3, 1/6]) * np.sqrt(3/m)
\end{lstlisting}\end{latin}
یک ماتریس تنک دستی با \(d=4,m=2,q=1\) داشته باشیم:
\[
 A=\begin{bmatrix}1&0&-1&1\\0&1&0&0\end{bmatrix},
 \quad x=\begin{bmatrix}2\\3\\1\\4\end{bmatrix},
 \quad Ax=\begin{bmatrix}5\\3\end{bmatrix}.
\]
هر ستون دقیقاً یک ناصفر دارد. خروجی اول \(2-1+4=5\) و خروجی دوم ۳ است. چون \(q=1\)، عامل \(1/\sqrt q\) این‌جا چیزی را تغییر نمی‌دهد.

اگر \(q=2\) باشد، هر ویژگی در دو خروجی متمایز با علامت تصادفی مشارکت می‌کند و \(P=A/\sqrt2\) است. برای تولید امن، seed به‌تنهایی یک عدد دلخواه آزمایش نیست؛ از transcript استخراج می‌شود. تولیدکننده پروتکل با SHAKE256، انتخاب‌های یکنواخت، رد برخورد سطرها و جداسازی دامنه کار می‌کند. درایه‌های یک ستون در سطرهای متمایز قرار می‌گیرند.

در مدار، از \(A\) صحیح استفاده می‌شود. اگر هم نقطه و هم مرکز با آن ضرب شوند، همه فاصله‌های مربعی \(q\) برابر فاصله‌های مربوط به \(P\) هستند؛ نزدیک‌ترین مرکز عوض نمی‌شود. برای گزارش مقدار هندسی نرمال‌شده یک‌بار بر \(q\) تقسیم می‌کنیم، نه دوبار.

\textbf{چرا جمع و تفریق جذاب است؟} برای محاسبه بومی، ضرایب صفر و علامت‌ها بسیاری از ضرب‌ها را حذف می‌کنند. اما در مدار فعلی، هزینه بررسی انتخاب سطر، احراز ماتریس و تعهد هم وجود دارد؛ نمی‌توان از همین مثال نتیجه گرفت کل اثبات حتماً به همان نسبت ارزان می‌شود.

\lesson{دانش صفر از صفر: چه چیزی را به چه کسی ثابت می‌کنیم؟}
اثبات دانش صفر یا ZK به‌طور شهودی یعنی «می‌توانم درستی یک ادعای دقیق را ثابت کنم، بدون این‌که اطلاعات خصوصی لازم برای محاسبه را به تو نشان بدهم». کسی که شاهد خصوصی را دارد \lr{Prover} و کسی که اثبات را بررسی می‌کند \lr{Verifier} است.

برای مثال فرض کن می‌خواهم ثابت کنم رمز یک قفل را می‌دانم، بدون گفتن رمز. یک نمایش امن باید طوری طراحی شود که فردی که رمز را ندارد نتواند با ترفند، verifier را قانع کند؛ و transcript نمایش هم رمز را لو ندهد. داستان قفل فقط شهود است و پروتکل رمزنگاری واقعی ما نیست.

سه ویژگی کلاسیک مهم‌اند: \textbf{کامل‌بودن} یعنی شاهد درست پذیرفته شود؛ \textbf{soundness} یعنی ادعای غلط با احتمال بیش از حد مجاز پذیرفته نشود؛ \textbf{دانش صفر} یعنی اثبات اطلاعات اضافه درباره شاهد، فراتر از خروجی عمومی تعریف‌شده، ندهد [۳]. تضمین‌های دقیق به مدل امنیتی و سیستم اثبات وابسته‌اند.

در پروژه ما، صورت ساده ادعا این است: «من همان داده‌ای را دارم که قبلاً به آن متعهد شده‌ام؛ این ماتریس تنک با تعهد عمومی سازگار است؛ داده و مرکزها را درست تبدیل کرده‌ام؛ انتساب‌ها نزدیک‌ترین مرکزهای تبدیل‌شده هستند و مجموع هزینه اعلام‌شده درست است.»

\textbf{شاهد خصوصی} شامل مختصات، salt و اطلاعات لازم برای محاسبه، مانند انتساب‌هاست. \textbf{اطلاعات عمومی} شامل تعهدها، مرکزها/پارامترهای عمومی سازگار با مدار و مقدار هزینه منتشرشده است. داده و برچسب‌های مثال آموزشی این فایل عمداً چاپ شده‌اند؛ این انتشار برای آموزش است، نه توصیه انتشار داده خصوصی در استفاده واقعی.

ZK صحت علمی منبع داده را اثبات نمی‌کند. اگر کسی از ابتدا داده ساختگی را commit کند، ممکن است تمام محاسبه روی همان داده ساختگی درست باشد. همچنین خروجی عمومی خود ممکن است اطلاعاتی بدهد؛ دانش صفر آن خروجی را پنهان نمی‌کند.

\takeaway{ZK درستی یک رابطه ریاضی را اثبات می‌کند؛ صداقت بیمارستان، کیفیت دیتاست، کاربرد پزشکی یا بهترین‌بودن خوشه‌ها را خودکار اثبات نمی‌کند.}

\lesson{یک مثال خیلی خودمانی برای دانش صفر: غار دو راهه}
یک غار را تصور کن که از ورودی به دو راه A و B تقسیم می‌شود. در عمق غار، یک درِ قفل‌شده این دو راه را به هم وصل می‌کند. من رمز در را می‌دانم و می‌خواهم این را به تو ثابت کنم، ولی نمی‌خواهم رمز را بگویم. تو بیرون می‌ایستی و نمی‌بینی اول وارد کدام راه می‌شوم.

من قبل از درخواست تو وارد یکی از راه‌ها می‌شوم. بعد تو \textbf{تصادفی} می‌گویی «از A برگرد» یا «از B برگرد». اگر رمز را بدانم، هرجا باشم می‌توانم در را باز کنم و از راهی که خواستی برگردم. اگر رمز را ندانم، فقط وقتی موفقم که درخواست تو اتفاقاً با راه اولیه‌ام یکی باشد؛ احتمال موفقیت حدود نصف است.

اگر این نمایش با انتخاب تازه چند بار تکرار شود، فرد بدون رمز باید هر بار خوش‌شانس باشد: ده بار موفقیت تصادفی احتمال \(2^{-10}\)، یعنی حدود یک از ۱۰۲۴ دارد. تو رمز را ندیدی؛ فقط شواهدی گرفتی که من امکان عبور از در را دارم.

\textbf{چرا ترتیب مهم است؟} اگر قبل از ورود به من بگویی از کدام راه برگردم، بدون رمز هم همان راه را انتخاب می‌کنم و نمایش ارزشی ندارد. این شهود به بخش post-commit نزدیک است: انتخاب اولیه باید پیش از چالش ثابت شده باشد.

این داستان یک \textbf{تشبیه آموزشی} است، نه شرح Groth16 و نه اثبات رسمی امنیت پروتکل ما. در سیستم واقعی به‌جای راه غار، رابطه ریاضی و شاهد داریم؛ به‌جای دیدن بازگشت، verifier یک proof و پارامترهای عمومی را بررسی می‌کند. همچنین صرف پنهان‌بودن رمز در داستان، جای تعریف رسمی شبیه‌سازی transcript در دانش صفر نیست.

\textbf{سه چیز متفاوت:} رمزنگاری معمولی پیام را طوری می‌پوشاند که دارنده کلید بتواند بازش کند؛ commitment مقدار را ثابت می‌کند؛ ZK درستی یک رابطه درباره مقدار خصوصی را اثبات می‌کند. hash، encryption و proof یک ابزار با سه اسم نیستند.

\takeaway{چالش تازه مانع تقلب آسان می‌شود؛ proof قرار نیست فایل مخفی را برای verifier رمزگشایی کند.}

\lesson{SNARK، Groth16 و «اثبات کوتاه» از دید یک تازه‌کار}
اسم‌ها ترسناک‌اند، ولی نقششان ساده‌تر است. مدار می‌گوید چه رابطه‌ای باید برقرار باشد. سیستم اثبات این رابطه و witness را به یک proof تبدیل می‌کند. verifier با اطلاعات عمومی و کلید بررسی می‌کند که proof معتبر است؛ لازم نیست همه داده خصوصی را داشته باشد و محاسبه اصلی را از نو انجام دهد.

در نام SNARK، ایده‌های «کوتاه‌بودن اثبات و بررسی»، «غیرتعاملی‌بودن در مدل موردنظر» و «استدلال دانستن شاهد» کنار هم‌اند. امنیت معمولاً محاسباتی است: prover با توان محاسباتی محدود و تحت فرض‌های رمزنگاری نباید ادعای غلط را اثبات کند. Groth16 یک سیستم مشخص در این خانواده است، نه اسم روش خوشه‌بندی ما. gnark ابزار Go برای نوشتن مدار و استفاده از این backendهاست [۳].

برای فهم هزینه، تهیه یک گواهی مهرشده را تصور کن: خواندن و بررسی تمام پرونده ممکن است گران باشد، اما گواهی نهایی یک برگ کوتاه است. proof ۱۹۶بایتی مثال کوچک هم همین هشدار را می‌دهد: کوتاه بودن سندِ نهایی، خواندن میلیون‌ها عدد و ساخت آن سند را ارزان نمی‌کند.

\textbf{Compile} شکل مدار را به سیستم قیود تبدیل می‌کند. \textbf{Setup} پارامترهای لازم سیستم اثبات را برای مدار یا دامنه سازگار می‌سازد. \textbf{Witness generation} مقادیر متغیرها را از داده آماده می‌کند. \textbf{Prove} اثبات را می‌سازد. \textbf{Verify} معتبر بودن آن را نسبت به ورودی عمومی بررسی می‌کند. برای ران علمی باید زمان هر مرحله و peak حافظه معلوم باشد.

اگر setup وابسته به مدار باشد، کلید یک شکل را نمی‌توان بدون بررسی برای شکل دیگر استفاده کرد. trusted setup هم یک فرض اعتماد دارد: اسرار مرحله راه‌اندازی نباید برای ساخت proof جعلی قابل استفاده باقی بمانند. در تست نرم‌افزار می‌توان پارامترهای آزمایشی ساخت؛ این کار اعتبار deploy تولیدی را ثابت نمی‌کند.

\takeaway{روش ما الگوریتم خوشه‌ای و پروتکل را طراحی می‌کند؛ Groth16 زیرساخت اثبات آن است. روش بومی سریع ممکن است مدار گران داشته باشد، و proof کوتاه ممکن است prover پرحافظه بخواهد.}

\lesson{تعهد، salt و درخت Merkle چه هستند؟}
تعهد را مثل گذاشتن داده در پاکتی تصور کن که بعداً نتوانی محتویاتش را عوض کنی. verifier فعلاً یک شناسه کوتاه از پاکت می‌بیند، نه کل داده را. در رمزنگاری به محدودیت تغییر محتوا \lr{binding} و پنهان‌ماندن محتوا \lr{hiding} می‌گوییم؛ تحقق دقیق این ویژگی‌ها به ساختار و فرض‌های رمزنگاری وابسته است.

یک hash ساده از داده‌ای با گزینه‌های کم ممکن است قابل حدس باشد. برای همین از salt تصادفی مخفی استفاده می‌شود. salt قرار نیست ماهیت محاسبه را عوض کند؛ از تطبیق مستقیم حدس‌ها با تعهد عمومی جلوگیری می‌کند.

در کد فعلی، مقادیر صحیح سطرها پشت هم قرار می‌گیرند. هر مقدار همراه salt وارد hash سازگار با میدان می‌شود و برگ درخت را می‌سازد. hash دو فرزند، والد را می‌سازد؛ این فرایند تا ریشه ادامه پیدا می‌کند:
\[
 h_i=\operatorname{MiMC}(x_i,\mathrm{salt}),\qquad
 h_{a,b}=\operatorname{MiMC}(h_a,h_b).
\]
برای تعداد غیرتوان‌دو، padding مشخص استفاده می‌شود. ریشه \(C_X\) به داده متعهد است. تغییر یک مقدار، مسیر hash تا ریشه را تغییر می‌دهد. مدار فعلی برای تطبیق ریشه، همه مقادیر و درخت لازم را پردازش می‌کند؛ این همان بخش پرهزینه است.

مرکزها با \(C_\theta\) و رکوردهای ماتریس با \(C_A\) احراز می‌شوند. hashهای تعهد در مسیر proof از نمایش میدان‌محور یکسان Python و Go استفاده می‌کنند. SHA256 در استخراج seed نقش جداگانه دارد؛ نباید آن را با hash داخل مدار یکی فرض کرد.

\textbf{دقت در توصیف نسخه فعلی:} پیش‌نویس برای \(\theta\) بسته بزرگ‌تری از پارامترها پیشنهاد می‌کند، اما تابع فعلی تعهد مرکزها را hash می‌کند؛ \(m,q\) به شکل مدار مقیدند و \(\mu\) غیرخالی در مسیر فعلی پشتیبانی نمی‌شود. ریشه داده نیز همه metadata پیشنهادی مقاله را به شکل یک بسته مستقل hash نمی‌کند. ادعای امنیت باید مطابق دامنه واقعی تعهد نوشته شود.

\takeaway{تعهد یک قرارداد برای ثابت‌ماندن داده است؛ ریشه کوتاه بودن به معنی ارزان بودن بازسازی و اثبات آن در مدار نیست.}

\lesson{مدار اثبات چرا با یک برنامه معمولی فرق دارد؟}
در Python می‌گوییم «این تابع را اجرا کن». در مدار، رابطه‌هایی می‌سازیم که witness باید همه آن‌ها را ارضا کند. مثلاً برای \(z=ax\) باید این رابطه و برای \(u=(z-c)^2\) رابطه مربع‌شدن برقرار باشد. برای nearest-centroid هم باید فاصله منتخب از فاصله مرکزهای دیگر بیشتر نباشد و قانون tie-break رعایت شود.

محاسبات مدار معمولاً در یک \textbf{میدان متناهی} انجام می‌شوند؛ یعنی عددها پیمانه یک عدد اول بزرگ هستند. برای شهود، در میدان کوچک با \(p=17\)، حاصل \(16+2\) برابر ۱ است، چون ۱۸ با کم‌کردن ۱۷ به ۱ تبدیل می‌شود. میدان واقعی پروژه BN254 است و بسیار بزرگ‌تر از ۱۷ است.

مشکل این است که فاصله و مقایسه را روی عددهای صحیح معمولی می‌خواهیم، نه عددهایی که می‌توانند دور پیمانه بپیچند. به همین دلیل محدودیت بازه لازم است: مختصات، اندیس‌ها، خروجی تبدیل، اختلاف‌ها و فاصله‌های مربعی باید در بازه‌ای باشند که معنی صحیح محاسبه حفظ شود.

در کد فعلی، مختصات ورودی و مرکزها باید \(|x_t|\le4096\) باشند. اندیس سطر باید غیرمنفی و کوچک‌تر از \(m\) باشد؛ عددی که در نمایش میدان معادل منفی یک است، اندیس معتبر نیست. برای اختلاف تصویرشده کران محافظه‌کارانه \(2dB\) و برای فاصله مربعی کران \(4md^2B^2\) در نظر گرفته می‌شود. جمع هزینه نیز با \(n\) رشد می‌کند. اصل موردنیاز، بزرگ‌تر بودن میدان نسبت به دو برابر کران محاسبات و مقایسه‌های مربوط است.

\textbf{چرا «جمع و تفریق ساده» کافی نیست؟} علامت، سطر انتخاب‌شده و ماتریس شاهد باید احراز شوند. نسخه فعلی routing تنک را با حلقه‌هایی روی \(n,m,d,q\) انجام می‌دهد؛ شمار عملیات منطقی آن را نباید با الگوریتم بومی \(O(ndq)\) یکی گرفت. کاهش هزینه مدار نیازمند اندازه‌گیری و بهینه‌سازی واقعی همین gadgetهاست.

\takeaway{صحت در میدان با صحت فاصله صحیح یکسان نیست، مگر این‌که محدودیت‌های بازه و نمایش داده درست طراحی شوند.}

\lesson{پنج مرحله پروتکل ما، به ترتیب دقیق}
ترتیب بخش امنیتی روش مهم است. اگر مالک داده بتواند بعد از دیدن ماتریس، داده را عوض کند، می‌تواند از ضعف‌های همان ماتریس سوءاستفاده کند.

\textbf{۱. Commit:} مالک داده، داده صحیح و salt را آماده می‌کند و \(C_X\) را منتشر می‌کند. مرکزها هم در \(C_\theta\) ثبت می‌شوند. هنوز چالش نهایی مشخص نشده است.

\textbf{۲. Challenge:} یک مقدار تازگی \(\rho\) از منبع تعریف‌شده می‌آید؛ مثلاً nonce از verifier. معتبر بودن این مرحله به قوانین واقعی منبع تصادف، احراز هویت و بودجه تلاش بستگی دارد.

\textbf{۳. Seed:} seed از زنجیره تعهد و چالش ساخته می‌شود:
\[
 s=\operatorname{SHA256}(\texttt{CommitJL-v1}\Vert C_X\Vert C_\theta\Vert\rho).
\]
تعهدهای میدان در این زنجیره با نمایش کانونی ۳۲بایتی وارد می‌شوند. تغییر داده متعهد، مرکز یا چالش، seed را تغییر می‌دهد.

\textbf{۴. Matrix authentication:} از seed ماتریس تنک تولید می‌شود؛ تعهد \(C_A\) و ساختار مجاز آن احراز می‌شود. همه طرف‌ها باید روی همان تولیدکننده، دامنه‌ها و encoding توافق داشته باشند.

\textbf{۵. Prove / Verify:} prover داده و salt و ماتریس و انتساب‌ها را به bridge می‌دهد؛ Go مدار را می‌سازد، شاهد را بررسی و اثبات تولید می‌کند؛ verifier proof را بررسی می‌کند. پیام نهایی باید به همان transcript متصل باشد.

\textbf{حد فعلی پیاده‌سازی:} استخراج seed و تطبیق ماتریس با آن در میزبان انجام می‌شود؛ مدار فعلی تولیدکننده XOF را از seed اجرا نمی‌کند. پس verifier کاربردی باید هم transcript و هم proof را بررسی کند. تأیید یک proof منفصل به‌تنهایی تمام زنجیره تازگی را اثبات نمی‌کند.

\takeaway{seed «عدد دلخواه آزمایش» نیست. زنجیره Commit تا Challenge و Matrix باید سالم بماند تا حمله انتخاب پس از دیدن ماتریس محدود شود.}

\lesson{چرا باید اول commit کنیم؟ یک حمله دوعددی}
فرض کن تبدیل از دو بُعد به یک بُعد برابر \(A=[1,-1]\) باشد. بردار \(v=(1,1)\) در فضای اصلی طول غیرصفر دارد، اما:
\[
 Av=1-1=0,\qquad \|v\|^2=2.
\]
اگر مالک داده اول ماتریس را ببیند، می‌تواند داده یا اختلاف مرکز را در جهتی انتخاب کند که تبدیل آن را نبیند. به این جهت «هسته ماتریس» می‌گوییم. همه تبدیل‌های خطی از فضای بزرگ‌تر به کوچک‌تر در حالت کلی چنین جهت‌هایی دارند.

اگر داده پیش از چالش متعهد شده باشد و چالش درست و غیرقابل پیش‌بینی تولید شود، مالک نمی‌تواند آزادانه پس از دیدن ماتریس، همان داده را جابه‌جا کند. این به معنی حذف تمام خطای کاهش بُعد نیست؛ فقط یک منبع مهم انتخاب تطبیقی را محدود می‌کند.

\textbf{حمله grinding:} ممکن است فرد چند بار تعهد یا چالش جدید بگیرد تا ماتریسی پیدا کند که اتفاقی بیشتر به نفعش است. اگر احتمال یک رویداد نامطلوب در هر تلاش \(p\) باشد، تحت استقلال، احتمال دیدن آن در \(Q\) تلاش:
\[
 1-(1-p)^Q.
\]
بدون نیاز به استقلال، اگر برای هر تلاش کران معتبر \(p\) داشته باشیم، union bound کران \(\min(1,Qp)\) می‌دهد. اما \textbf{نرخ تجربی خطای فاصله در ۲۵۶ ماتریس، همان احتمال پذیرش proof جعلی یا کران رمزنگاری نیست}.

برای محدودکردن تلاش، TaskLog محلی زمان و بودجه را ثبت می‌کند. این ابزار فعلی جای یک سامانه پایدار، احراز هویت‌شده و مقاوم در برابر چند هویت یا چند فرایند را نمی‌گیرد. deploy واقعی به سیاست جلوگیری از retry نامجاز نیاز دارد.

\takeaway{post-commit امنیت ترتیب انتخاب را بهتر می‌کند؛ هیچ تعهدی تضمین نمی‌کند هر projection روی هر دیتاست هندسه را خوب حفظ کند.}

\lesson{nonce، beacon و commit-reveal چه تفاوتی دارند؟}
\textbf{Nonce:} verifier بعد از commit یک عدد تازه می‌فرستد. ساده است، ولی باید به تازگی و رفتار verifier یا سیاست پروتکل اعتماد کنیم. تکرار یا انتخاب جانبدارانه باید در مدل تهدید دیده شود.

\textbf{Beacon:} به خروجی یک منبع عمومیِ زمان‌مند ارجاع می‌دهیم؛ مثلاً round مشخص یک سرویس تصادف عمومی. انتخاب round، دسترس‌پذیری و زمان انتشار اهمیت دارند. \textbf{در آزمایش ما beacon یک trace محلی شبیه‌سازی‌شده است}. تأخیر شبکه و امنیت سرویس beacon واقعی اندازه‌گیری نشده است؛ عنوان لازم در گزارش \lr{Simulated local trace} است.

\textbf{Commit-reveal:} چند طرف ابتدا به سهم‌های تصادف خود متعهد می‌شوند؛ سپس سهم‌ها را reveal می‌کنند و با هم ترکیب می‌شوند. اگر سهم مناسبِ یک طرف صادق تا موعد پنهان بماند، می‌توان به تولید تصادف کمک کرد. اما طرف آخر ممکن است بعد از دیدن سهم‌ها از reveal صرف‌نظر کند؛ مشکل abort و دسترس‌پذیری باید جدا حل شود.

پیاده‌سازی آزمایش commit-reveal منطق تعهد و آشکارسازی را جدا از nonce طی می‌کند، ولی همه شرکت‌کنندگان در محیط محلی شبیه‌سازی شده‌اند. این یک شبکه واقعی از نهادهای مستقل نیست.

این سه روش مسیرهای یکسان با اسم‌های متفاوت نیستند؛ سیاست تولید \(\rho\) فرق دارد. با این حال انتظار نداریم تغییر منبع تصادف معتبر، خودش دقت هندسی را به‌صورت جادویی افزایش دهد. سؤال اصلی RQ3 هزینه اضافه و اثر تلاش‌های دوباره است، نه این‌که beacon از PCA دقیق‌تر باشد.

\takeaway{نام یک منبع تصادف، به‌تنهایی گواه امنیت آن نیست؛ نحوه اجرا و برچسب شبیه‌سازی باید مشخص باشد.}

\lesson{JL دقیقاً چه قولی می‌دهد؟}
ایده Johnson--Lindenstrauss این است که برای یک مجموعه متناهی از نقطه‌ها، یک تبدیل مناسب با احتمال بالا می‌تواند فاصله‌ها را تقریباً نگه دارد. روی مجموعه بردارهای موردنیاز فرض می‌کنیم:
\[
 (1-\varepsilon)\|v\|^2\le\|Pv\|^2
 \le(1+\varepsilon)\|v\|^2.
\]
\(\varepsilon=0.15\) یعنی فاصله \textbf{مربعی} بین ۸۵ و ۱۱۵ درصد مقدار اصلی باشد. این بازه را نباید مستقیم به فاصله غیرمربعی نسبت داد؛ آن‌جا ریشه عوامل مطرح است.

در نظریه کلاسیک، مرتبه بُعد موردنیاز به \(\log N/\varepsilon^2\) وابسته است؛ \(N\) اندازه مجموعه مرتبط است. جزئیات احتمال، ثابت‌ها و ساختار تنکی مهم‌اند [۸]. این قانون کلی به ما اجازه نمی‌دهد هر \(m,q\) کوچک را روی هر داده‌ای «تضمین‌شده» بنامیم.

برای ماتریس تنک فعلی، نرمال‌سازی و علامت‌های مناسب باعث می‌شوند امید ریاضیِ طول مربعی حفظ شود. ولی حفظ \textbf{امید ریاضی} با تمرکز احتمال و حفظ هم‌زمان همه فاصله‌ها فرق دارد. ممکن است میانگین چند draw خوب باشد و یک draw برای بردار خاص بد عمل کند.

اثبات رابطه JL و استفاده از قضیه آن برای الگوریتم خوشه‌ای باید دو مرحله داشته باشد: اول، آیا رویداد هندسی روی مجموعه لازم برقرار است؟ دوم، اگر برقرار بود چه نتیجه‌ای درباره تصمیم‌ها داریم؟ مدار فعلی همه رویدادهای فاصله اصلی و تصویرشده را به‌عنوان یک تضمین جهانی داخل proof احراز نمی‌کند.

\textbf{هزینه ثابت‌مرکز:} اگر همین کران برای تمام \(x_i-c_j\) برقرار باشد، گرفتن کمینه و جمع‌زدن، هزینه نرمال‌شده ثابت‌مرکز را هم در بازه \((1\pm\varepsilon)\) هزینه اصلی نگه می‌دارد. این قضیه درباره همان مرکزهاست؛ بهینه‌بودن آموزش جدید را ثابت نمی‌کند.

\takeaway{قضیه «اگر رویداد هندسی برقرار باشد...» با ادعای «روی همه ران‌های ما آن رویداد برقرار است» فرق دارد.}

\lesson{گواه پایداری انتساب: چه موقع گروه عوض نمی‌شود؟}
فرض کن برای یک نقطه فاصله مربعی تا نزدیک‌ترین مرکز \(D_1\) و تا دومین مرکز \(D_2\) باشد. هرچه فاصله دوم نسبت به اول بزرگ‌تر باشد، تصمیم مقاوم‌تر است. اگر projection خطای نسبی حداکثر \(\varepsilon\) داشته باشد، شرط کافی این است:
\[
 (1+\varepsilon)D_1<(1-\varepsilon)D_2,
 \qquad \frac{D_2}{D_1}>\frac{1+\varepsilon}{1-\varepsilon}.
\]
برای \(\varepsilon=0.15\)، نسبت لازم حدود ۱٫۳۵۳ است. اگر \(D_1=4,D_2=9\)، حتی بدترین افزایش اول و کاهش دوم به ۴٫۶ و ۷٫۶۵ می‌رسند، پس مرکز اول همچنان نزدیک‌تر است.

اگر \(D_1=4,D_2=4.4\) باشد، شرط کافی برقرار نیست. این \textbf{نمی‌گوید حتماً انتساب عوض می‌شود}؛ فقط این گواه برای تضمین آن نقطه کافی نیست. نقطه‌های نزدیک مرز تصمیم، سخت‌تر گواه می‌گیرند. حالت فاصله صفر یا tie نیز باید با قواعد مقایسه دقیق پیاده‌سازی بررسی شود، نه تقسیم کورکورانه بر صفر.

\textbf{پوشش گواه} یعنی چند درصد نقطه‌ها هم شرایط هندسی موردنیاز و هم شرط margin را دارند. \textbf{صحت روی مجموعه eligible} یعنی میان همان نقطه‌های واجد شرایط چند تغییر دیده شده است. ممکن است صحت ۱۰۰ درصد باشد ولی فقط یک درصد نقطه‌ها eligible باشند؛ بنابراین باید هر دو عدد را گزارش کنیم.

در RQ4 بررسی تجربی ما، روی مجموعه‌های eligible تغییر مشاهده نشده است. اما این نتیجه به معنی «تمام دیتاست گواه دارد» نیست. همچنین زیادکردن \(\varepsilon\) از یک طرف احراز خطای projection را آسان‌تر و از طرف دیگر شرط margin را سخت‌تر می‌کند؛ پوشش لزوماً صعودی نیست.

\takeaway{گواه پایداری یک شرط کافی برای نقطه‌های دارای حاشیه امن است؛ درباره نقطه‌های خارج آن، ممکن است جواب درست باشد ولی تضمین نداریم.}

\lesson{DBSCAN: خوشه با همسایه‌ها، نه با مرکزها}
در DBSCAN یک شعاع \(\rho\) و حداقل تعداد همسایه \lr{minPts} داریم. اگر اطراف نقطه به اندازه کافی نمونه باشد، نقطه هسته‌ای است؛ اتصال نواحی متراکم خوشه می‌سازد و بعضی نقطه‌ها نویز باقی می‌مانند. برخلاف KMeans لازم نیست از ابتدا تعداد خوشه‌ها را مشخص کنیم [۱۲].

تصمیم پایه آن این است: «فاصله دو نقطه از \(\rho\) کمتر است یا نه؟» با \(\rho=1\)، نقطه‌های با فاصله مربعی ۰٫۸ همسایه و با ۱٫۳ غیرهمسایه‌اند. فاصله مربعی ۱٫۰۰۱ به مرز بسیار نزدیک است و یک تغییر کوچک آن را جابه‌جا می‌کند.

اگر کران JL با \(\varepsilon=0.15\) برقرار باشد، ناحیه خاکستری فاصله \textbf{مربعی} برای همسایگی:
\[
 \left[\frac{\rho^2}{1+\varepsilon},\frac{\rho^2}{1-\varepsilon}\right]
 =[0.869565,1.176471].
\]
پایین این ناحیه همسایه‌بودن و بالای آن غیرهمسایه‌بودن قابل گواه است؛ داخل آن تضمین حاصل نمی‌شود. برای فاصله غیرمربعی حدود ۰٫۹۳۳ تا ۱٫۰۸۵ مطرح است.

اگر تمام روابط همسایگیِ لازم حفظ شوند، می‌توان با شرایط و قراردادهای الگوریتم درباره حفظ خروجی DBSCAN استدلال کرد. اما \textbf{آزمایش فعلی RQ4 روابط جفتی را بررسی می‌کند، نه یک مدار اثبات کامل DBSCAN}. نمی‌توان جدول پایداری جفت‌ها را «اثبات کامل خوشه‌بندی DBSCAN» نامید.

دیتاست مصنوعی موجود عمداً ۲۵۰ جفت نزدیک مرز در میان ۱۰۰۰ نقطه می‌سازد. تعداد کل جفت‌ها ۴۹۹۵۰۰ است؛ پس جفت‌های سخت فقط حدود ۰٫۰۵ درصد کل جفت‌ها هستند. میانگین عالی روی همه جفت‌ها ممکن است ضعف روی همین جفت‌های مهم را پنهان کند. به همین دلیل نتیجه ناحیه مرزی را جدا نگاه می‌کنیم.

\takeaway{برای DBSCAN، حفظ همسایگی نزدیک مرز مهم‌تر از یک درصد کلی است که عمدتاً از جفت‌های دور تشکیل شده است.}

\lesson{اعشار، quantization و مشکل عددهای کوچک}
مدارهای این پروژه از نمایش صحیح استفاده می‌کنند. برای تبدیل اعشار به صحیح، اول مقیاس می‌دهیم و بعد گرد می‌کنیم:
\[
 x_{\mathrm{int}}=\operatorname{round}(Sx).
\]
اگر فاصله بین ۰٫۹ و ۰٫۱ را بخواهیم، مقدار مربعی درست \((0.9-0.1)^2=0.64\) است. تبدیل زودهنگام هر دو به int، آن‌ها را صفر می‌کند و جواب اشتباه صفر می‌دهد. با \(S=10\)، دو مقدار ۹ و ۱ می‌شوند؛ مربع اختلاف ۶۴ است و با تقسیم بر \(S^2=100\)، مقدار ۰٫۶۴ برمی‌گردد.

در کد فعلی باید مقیاس‌دهی و rounding پیش از تبدیل انجام شود. مثال ۱۰ صرفاً برای فهم است؛ برخی helperهای رسمی مقیاس‌های توان‌دو را می‌پذیرند. هر نمونه و مرکز باید با یک قرارداد مقیاس سازگار باشد. اگر ماتریس هم \(A\) به‌جای \(P=A/\sqrt q\) است، هزینه حقیقی از مقدار صحیح با عامل‌های مقیاس و \(q\) مناسب بازسازی می‌شود.

\textbf{Overflow در Python و NumPy:} int معمولی Python دقت دلخواه دارد، ولی \lr{numpy.int64} محدود است. مربع و جمع روی بردار بزرگ می‌تواند دور بپیچد و حتی فاصله منفی یا انتساب غلط تولید کند. پیاده‌سازی فاصله با کران امن مسیر سریع int64 را انتخاب می‌کند و در صورت نیاز از محاسبه صحیحِ بدون overflow استفاده می‌کند. اتکا به \lr{numpy.int128} راه‌حل قابل‌حمل Windows نیست.

در مدار، مسئله جداست: حتی اگر Python overflow نکند، باید wrap میدان هم با range check مهار شود. آزمایش مرزهای \(B=4096\) و بعدهای بزرگ و تست شاهد نامعتبر، هرکدام لایه متفاوتی از صحت را بررسی می‌کنند.

\textbf{دقت بیت در ران ما:} با \(B\) ثابت، تغییر هم‌زمان مقیاس فیزیکی و \(S\) می‌تواند همان witness صحیح را بسازد. پس برچسب «۸ بیت» همیشه به معنی مختصات واقعاً ۸بیتی یا قید کمتر نیست. این نکته برای تفسیر RQ5 حیاتی است.

\lesson{یک proof واقعی کوچک که برای این راهنما اجرا شد}
برای این فایل یک نمونه چهارردیفی با چهار ویژگی، دو مرکز و تبدیل به دو بُعد اجرا کردیم. این تنها یک مثال آموزشی و عمومی است؛ داده‌های آن عمداً قابل مشاهده‌اند.
\codefile{proof_output.txt}
زنجیره Commit، Challenge، Seed، Matrix Auth و تولید و تأیید proof موفق بود. خروجی واقعی شامل \textbf{\lr{\ProofConstraints} قید} و proof با \textbf{\lr{\ProofBytes} بایت} برای backend \lr{Groth16} بود. زمان دیواری این یک اجرا حدود \lr{\ProofWall} ثانیه شد. اندازه کوچک proof به معنی کوچک بودن حافظه ساخت آن نیست؛ prover باید مدار و witness را پردازش کند.

برای همین ماتریسِ این اجرا، نقطه‌ها و مرکزهای تصویرشده و هزینه‌های منتخب به شکل زیر هستند. این جمع را می‌توان با دست بررسی کرد. proof همین محاسبه را همراه ارتباط با تعهدها و بازه‌های مجاز احراز می‌کند؛ هزاران قید صرفاً برای چهار جمع ساده نیست.
\input{proof_math.tex}

\begin{latin}\begin{lstlisting}
import hashlib
from commitjl.protocol import run_protocol
result = run_protocol(Xsmall, Csmall, m=2, q=1, scale=1,
    generate_zk_proof=True, proof_timeout=180,
    challenge_fn=lambda: hashlib.sha256(
        b"educational-guide-challenge").digest())
assert result["proof_result"].verify_ok
\end{lstlisting}\end{latin}
در این قطعه، \lr{Xsmall} و \lr{Csmall} همان ماتریس‌های چاپ‌شده بالا هستند. برای تعریف کامل و قابل اجرای آن‌ها از اسکریپت نمونه استفاده می‌کنیم.

\clearpage\subsection{از این proof چه می‌فهمیم و چه نمی‌فهمیم؟}
چهار نمونه بالا عمداً ساده‌اند تا بتوانیم از ابتدا تا انتها محاسبه را با دست دنبال کنیم. مدار با خروجی میزبان ارتباط دارد: ریشه داده باید با همان \(X\) و salt سازگار باشد، مرکزها با تعهد مرکزها سازگار باشند و ماتریس هم ساختار مجاز و تعهد درست داشته باشد. بعد باید مختصات تصویرشده، انتساب‌ها و جمع هزینه با این ورودی‌ها جور باشند. صرفاً این‌که یک نفر عدد ۲ را تایپ کند چنین ارتباطی را ثابت نمی‌کند.

اگر ادعا کنیم نمونه دوم به مرکز دوم نزدیک‌تر است، باید همه فاصله‌ها و مقایسه مربوط آن ادعا را پشتیبانی کنند. یک witness با محاسبه غلط نباید solver را راضی کند. اگر یک مختصه خارج بازه وارد کنیم، مجاز نیست فقط آن را به پیمانه میدان تبدیل کنیم و بگوییم «یک عدد میدان شد»؛ range check برای منع همین رفتار است. این‌ها توضیح رابطه‌ای هستند که مدار باید enforce کند، نه نتیجه تست حمله‌ای جدید در این فایل.

در استفاده خصوصی، verifier نیازی ندارد همین جدول چهارردیفی را ببیند. جدول این‌جا برای آموزش چاپ شده است تا بفهمیم witness چیست. verifier واقعی ورودی‌های عمومی و proof را می‌بیند؛ همچنین باید زنجیره transcript را در میزبان بررسی کند. این موضوع با محرمانگی خروجی هزینه فرق دارد: اگر هزینه عمومی باشد، هر اطلاعاتی که خود آن می‌دهد جزو خروجی مجاز است.

این اجرا نه ده تکرار بنچمارک است و نه مقایسه SparseJL و PCA در RQ2. چالش آن برای آموزش ثابت انتخاب شده؛ سیاست امن تازگی در deploy واقعی را از آن نتیجه نمی‌گیریم. salt تصادفی است، بنابراین ریشه‌ها، seed و ماتریس ممکن است در تکرار بعد عوض شوند.

در Groth16 مراحل compile، setup، ساخت witness، prove و verify جدا هستند. setup باید برای مدار سازگار مدیریت شود؛ استفاده دوباره برای همان شکل می‌تواند هزینه را کم کند. trusted setup، مفروضات اعتماد و مدیریت کلید می‌خواهد. backend آزمایشی PLONK در این مخزن از SRS ناامنِ تست استفاده می‌کند و مبنای deploy امن نیست.

\takeaway{اکنون یک اتصال واقعی transcript به proof در نمونه کوچک داریم، اما موفقیت آن گلوگاه مدار صدها هزار مختصه‌ای را حل نمی‌کند.}

\lesson{فهرست دیتاست‌ها و حدود دقیق ران فعلی}
این جدول مربوط به \textbf{آخرین ران تشخیصی محلی} است، نه وعده پوشش کامل تمام اندازه‌های مقاله. نمونه آموزشی همین فایل هم با این اندازه‌ها بارگذاری شده است.
\begin{latin}\begin{center}\small
\begin{tabular}{lrrl}
\toprule Dataset & Local \(n\) & \(d\) & Source / pool \\\midrule
Gaussian mixture & 1,000 & 256 & Generated instance \\
DBSCAN synthetic & 1,000 & 128 & Generated boundary pairs \\
MNIST & 1,000 & 784 & 70,000 images \\
Fashion-MNIST & 1,000 & 784 & 70,000 images \\
Gisette & 1,000 & 5,000 & 6,000 loaded training rows \\
Arcene & 100 & 10,000 & 100 loaded training rows \\
20 Newsgroups TF-IDF & 1,000 & 2,000 & 11,314 training documents \\
\bottomrule\end{tabular}\end{center}\end{latin}
منبع Fashion و MNIST معمولاً ۶۰هزار train و ۱۰هزار test دارد. بارگذار فعلی برای این بررسی از pool موجود نمونه می‌گیرد؛ این ارزیابی طبقه‌بندی با test جدا نیست. Gisette در مستندات منبع مجموع ۷هزار ردیف train/validation برچسب‌دار دارد، اما مسیر فعلی فقط training شش‌هزاری را بارگذاری می‌کند. Arcene در منبع مجموع بزرگ‌تری تا ۹۰۰ نمونه با بخش‌های مختلف دارد؛ مسیر فعلی training صدتایی را دارد، نه همه ۹۰۰ نمونه.

در متن، \(n=100\) Arcene «صد نمونه از pool صدتایی بارگذاری‌شده» است؛ آن را full کل دیتاست مرجع نمی‌نامیم. در متن TF-IDF، تعداد ویژگی به واژگان انتخاب‌شده وابسته است و مثل ۷۸۴ پیکسل یک ویژگی ذاتی ثابت نیست.

\textbf{Native چیست؟} محاسبه با Python/NumPy روی داده، بدون ساخت proof. Native می‌تواند روی دیتاست کامل یا زیرمجموعه اجرا شود. در این ران، native روی اندازه‌های جدول اجرا شده است. «native» مساوی «sample» نیست؛ اندازه تنظیم‌شده تعیین می‌کند sample هست یا نه.

\textbf{چه چیزهایی اجرا شدند؟} RQ1، RQ3، RQ4 و RQ5 این کمپین native هستند. این کمپین هیچ ردیف RQ2 جدیدی اجرا نکرده است. تلاش‌های بزرگ RQ2 قبلی ناموفق بوده‌اند و نتیجه موفق proof بزرگ نداریم.

\lesson{داده Gaussian و DBSCAN مصنوعی چه شکلی دارند؟}
در Gaussian mixture خودمان چند مرکز تصادفی می‌سازیم و اطراف هرکدام نویز نرمال اضافه می‌کنیم. هر سطر یک بردار مصنوعی است؛ مثلاً یک مشتری یا تصویر واقعی نیست. label می‌گوید مولد این نقطه را از کدام مرکز ساخته است، ولی خوشه‌بندی هنگام اجرا آن label را نمی‌بیند.
\codefile{gaussian_mixture_sample.txt}
این دوازده مقدار ابتدایی اولین بردار واقعی تولیدشده در مثال‌اند؛ ۲۴۴ ویژگی دیگر چاپ نشده‌اند. چون داده مصنوعی است، می‌توانیم فاصله مرکزها و نویز را کنترل کنیم و بفهمیم روش در مسئله آسان یا سخت چطور رفتار می‌کند.

پارامتر \(\gamma\) در مولد جدایی مراکز را افزایش می‌دهد. عدد کوچک یعنی گروه‌ها هم‌پوشانی بیشتری دارند؛ عدد بزرگ یعنی معمولاً آسان‌تر از هم جدا می‌شوند. این پارامتر یک کران قطعی حداقل فاصله همه مرکزها نیست. seed هم خروجی تکرارپذیر مولد را تعیین می‌کند.

دیتاست DBSCAN مصنوعی برای آزمون مرز همسایگی طراحی شده است: تعدادی نقطه دور از هم و تعدادی جفت با فاصله نزدیک شعاع تصمیم. این دیتا مسئله دیگری را می‌سنجد، حتی اگر شکل ذخیره‌اش باز هم ماتریس عددی باشد.
\codefile{dbscan_margin_sample.txt}
برای بررسی دقیق generator، پارامتر، نسخه و hash باید با هم ثبت شوند. برابر بودن seed بدون یکسان بودن ترتیب فراخوانی‌های تصادفی و پیاده‌سازی، hash یکسان را تضمین نمی‌کند. نمونه‌ها و hashهای منبع در فایل \lr{examples.json} ذخیره شده‌اند.

\takeaway{داده مصنوعی برای جداکردن اثر پارامترها عالی است، ولی موفقیت روی آن جای نتیجه روی داده واقعی را نمی‌گیرد.}

\lesson{Fashion-MNIST: لباس‌ها، نه رقم‌ها}
Fashion-MNIST از نظر شکل ذخیره شبیه MNIST است: تصویر خاکستری \(28\times28\)، یک بردار ۷۸۴پیکسلی، و برچسب یکی از ده دسته لباس. دسته‌ها شامل تی‌شرت، شلوار، پلیور، پیراهن، کت، صندل، پیراهن مردانه، کتانی، کیف و بوت هستند [۵].

\chart{fashion_mnist_samples.png}{تصویرهای واقعی بارگذاری‌شده؛ عدد عنوان، کد کلاس لباس است.}
\codefile{fashion_mnist_sample.txt}
بسیاری از مؤلفه‌های ابتدا صفرند، چون گوشه تصویر پس‌زمینه است. این نشانه ناقص بودن دانلود نیست. برای دیدن ساختار، باید بردار را به \(28\times28\) برگردانیم. همان کد \lr{reshape} که برای رقم استفاده شد این‌جا هم مناسب است.

در آزمایش، نقطه‌های شبیه از نظر پیکسل به مرکزهای KMeans نزدیک می‌شوند. این دسته‌بندی شاید تا حدی با دسته لباس مطابق باشد، ولی تضمین نمی‌کند پنج خوشه فعلی دقیقاً پنج گروه معنایی مشخص باشند. تی‌شرت و پیراهن مردانه از نظر تصویر می‌توانند نزدیک باشند، در حالی که label متفاوت دارند.

روی این داده در ران اخیر، SparseJL در \(m=64,q=2\) توافق متوسط ARI حدود ۰٫۸۴۸ با مرجع اصلی داشت. PCA در همان بُعد حدود ۰٫۹۹۹ داشت. پس این یک منطقه کاربردی برای روش ماست، اما ادعای «از PCA دقیق‌تر است» از این نتایج درنمی‌آید.

\takeaway{پیکسل‌های تصویر اطلاعات زیادی دارند، اما خوشه‌بندی پیکسلی با درک مفهوم لباس یا آموزش classifier یکسان نیست.}

\lesson{Gisette و Arcene: بردارهای خیلی بلند}
\textbf{Gisette} یک benchmark انتخاب ویژگی مربوط به تفکیک رقم‌های دست‌نویس ۴ و ۹ است. بردار هر نمونه ۵۰۰۰ ویژگی دارد؛ این‌ها فقط تصویر خام \(28\times28\) نیستند. ویژگی‌های پردازش‌شده و ویژگی‌های probe نامربوط در داده وجود دارند [۴]. هدف منبع، طبقه‌بندی دودویی است، اما ما روی بردارها خوشه‌بندی و حفظ تصمیم را بررسی می‌کنیم.
\codefile{gisette_sample.txt}
وجود تنها دو برچسب اصلی، ما را مجبور نکرده در ران فعلی \(k=2\) بگذاریم؛ \(k=5\) در بسیاری از این مقایسه‌ها برای تعریف مشترک آزمایش استفاده شده است. در نتیجه پنج خوشه را نباید پنج نوع رقم نامید. SparseJL در \(m=64\) برای این داده توافق حدود ۰٫۴۶۵ داشت و PCA حدود ۰٫۹۵۹؛ یک محدودیت واقعی روش در این تنظیم است.

\textbf{Arcene} داده برداری طیف‌سنجی جرمی با ۱۰هزار ویژگی و مسئله مرجع تمایز سرطان/غیرسرطان است؛ ویژگی‌های probe هم دارد [۶]. هر ردیف یک تصویر بیماری نیست؛ مجموعه‌ای از شدت‌های عددی ویژگی‌های طیف و ستون‌های افزوده‌شده است.
\codefile{arcene_sample.txt}
Arcene در حفظ خوشه‌بندی قبل و بعد projection خوب ظاهر شده، اما نتیجه پزشکی نداریم. ARI مرجع KMeans نسبت به برچسب واقعی در این نمونه حدود ۰٫۰۳۲ است. پس ARI نزدیک یک بعد از projection عمدتاً می‌گوید «همان تصمیم‌های KMeans را نگه داشتیم»، نه «سرطان را با دقت بالا تشخیص دادیم».

\takeaway{بالابودن \(d\) به‌تنهایی موفقیت SparseJL را تضمین نمی‌کند: Arcene و Gisette هر دو پُربُعدند ولی نتایج بسیار متفاوتی دارند.}

\lesson{متن به بردار: 20 Newsgroups و TF-IDF}
در 20 Newsgroups هر نمونه یک نوشته درباره یکی از بیست موضوع است. متن ابتدا به عدد تبدیل می‌شود. \lr{TF-IDF} به واژه‌ای که در یک نوشته مهم و در همه نوشته‌ها خیلی رایج نیست وزن بیشتری می‌دهد. در بارگذار ما header، footer و quote حذف می‌شوند، stopwordهای انگلیسی کنار گذاشته می‌شوند و حداکثر ۲۰۰۰ ویژگی واژگانی نگه داشته می‌شود [۷].

برای فهم فرمت، این چهار متن \textbf{ساختگی آموزشی} را تبدیل کرده‌ایم؛ این جمله‌ها نقل از دیتاست اصلی نیستند:
\begin{latin}\begin{lstlisting}
from sklearn.feature_extraction.text import TfidfVectorizer
docs = ["car engine wheel", "car wheel road",
        "space rocket orbit", "space rocket moon"]
v = TfidfVectorizer()
X_text = v.fit_transform(docs).toarray()
print(v.get_feature_names_out())
print(X_text.round(3))
\end{lstlisting}\end{latin}
\codefile{tfidf_output.txt}
هر ستون یک واژه است؛ سطر خودرو در ستون‌های space و rocket صفر دارد. سطرهای این مثال طول نرمال‌شده یک دارند. در متن واقعی هم ماتریس بسیار تنک است، چون هر نوشته تنها تعداد کمی از واژگان کل را دارد.

\textbf{بردار واقعی از زیرمجموعه ما:}
\codefile{20ng_tfidf_sample.txt}
واژگان در مسیر فعلی روی pool training یاد گرفته می‌شوند و بعد هزار سطر نمونه گرفته می‌شود؛ این یک آزمون پیش‌بینیِ held-out مستقل نیست. ماتریس برای بخش‌هایی از pipeline به نمایش متراکم تبدیل می‌شود؛ تنک بودن محتوایی الزاماً به حافظه تنک منجر نشده است.

در \(m=64\)، توافق SparseJL حدود ۰٫۲۲۳ و PCA حدود ۰٫۶۹۰ است. این اعداد درباره پنج مرکز مرجع هستند، نه تشخیص بیست موضوع. برای ادعای کاربرد روی embeddingهای زبانی جدید به آزمایش جدا نیاز داریم؛ چنین دیتاستی در این کمپین اجرا نشده است.

\lesson{معیارها: ARI، distortion، cost ratio و زمان}
\textbf{ARI} شباهت دو تقسیم‌بندی را با درنظرگرفتن جفت‌نقطه‌ها و اصلاح توافق تصادفی اندازه می‌گیرد. ARI برابر ۱ یعنی تقسیم‌بندی‌های یکسان تا جابه‌جایی نام خوشه‌ها. مقدار نزدیک صفر یعنی توافق در حد مرجع تصادفیِ تعریف معیار؛ مقدار منفی هم ممکن است. این عدد درصد نمونه‌های درست نیست.

در RQ1 معیار اصلی \(\operatorname{ARI}(a_{\rm full},a_{\rm projected})\) است. ARI نسبت به label واقعی یک ستون ارزیابی دیگر است. اگر مرجع اصلی خودش نسبت به واقعیت ضعیف باشد، حفظ آن مرجع الزاماً فایده کاربردی بالایی ندارد.

\textbf{Distortion} تغییر نسبی فاصله را می‌سنجد:
\[
 \delta(v)=\left|\frac{\|Pv\|^2}{\|v\|^2}-1\right|.
\]
بردار صفر باید جدا مدیریت شود. صدک ۹۵ درصدیِ distortion می‌گوید ۹۵ درصد بردارهای بررسی‌شده از آن مقدار کمترند، نه این‌که حداکثر خطا یا تضمین جهانی همان است. اگر برای نمونه‌ای delta حدود ۰٫۳۳ باشد، شرط \(\varepsilon=0.15\) برای همه آن فاصله‌ها برقرار نیست.

\textbf{Cost ratio} باید با تعریف دقیق گزارش شود. در pipeline جدید، نسبت objective پس از ساخت مرکزها از برچسب‌های تصویرشده و ارزیابی آن مرکزها در فضای اصلی، به objective مرجع سنجیده می‌شود. این دقیقاً همان \(\Psi/q\) برای مرکزهای ثابت در قضیه نیست. مقدار نزدیک یک می‌تواند با ARI پایین همراه باشد؛ دو تقسیم‌بندی متفاوت ممکن است هزینه تقریباً مساوی داشته باشند.

\textbf{زمان} هم چند نوع دارد: آموزش مرجع، fit PCA، تولید ماتریس، projection، assignment، commitment، compile، setup، prove و verify. زمان assignment کوچک به‌تنهایی کل هزینه پروتکل یا فایده ZK را نشان نمی‌دهد. \textbf{حافظه peak} یعنی بیشترین حافظه در اجرا، نه اندازه فایل proof.

\takeaway{باید همیشه بپرسیم «این عدد نسبت به چه مرجعی، روی چه مجموعه‌ای و برای کدام مرحله محاسبه شده؟»}

\lesson{پنج سؤال پژوهشی مقاله، به زبان ساده}
\textbf{RQ1 ـ کیفیت و کاهش بُعد:} اگر بردار را کوتاه کنیم، تصمیم‌های نزدیک‌ترین مرکز چقدر حفظ می‌شوند؟ SparseJL در برابر PCA و baselineها کجا جواب قابل قبول می‌دهد؟ چه \(m,q\)هایی کافی یا ناکافی‌اند؟ کد با مرکزهای مرجع مشترک ARI، objective و distortion را مقایسه می‌کند.

\textbf{RQ2 ـ هزینه اثبات:} برای همان ادعا و همان ورودی‌ها، ساخت و بررسی proof چه زمان، تعداد قید و حافظه‌ای می‌خواهد؟ کاهش فاصله‌ها از \(d\) به \(m\) آیا واقعاً در کل پروتکل اثر دارد؟ این سؤال مزیت اصلی \textbf{اثباتی} را می‌سنجد؛ سؤال اصلی دقت نیست. باید full، SparseJL و PCA با تعهد و دامنه ادعای منصفانه مقایسه شوند.

\textbf{RQ3 ـ امنیت ترتیب و سربار:} Commit قبل از چالش چه هزینه‌ای اضافه می‌کند؟ منابع تصادف چه تفاوت عملی دارند؟ با تلاش دوباره چه‌قدر احتمال یافتن projection نامطلوب بالا می‌رود؟ کد زمان مراحل میزبان و رویدادهای distortion در drawهای مستقل را می‌سنجد؛ حمله شکستن رمزنگاری انجام نمی‌دهد.

\textbf{RQ4 ـ گواه پایداری:} در چه شرایطی می‌توانیم بگوییم انتخاب مرکز یا همسایه تغییر نمی‌کند؟ چند درصد نمونه‌ها از این گواه بهره می‌برند؟ کد شرط خطای هندسی و حاشیه تصمیم را با تغییرات واقعی مقایسه می‌کند؛ برای DBSCAN جفت‌ها و مرز شعاع را بررسی می‌کند.

\textbf{RQ5 ـ حساسیت پارامترها:} جدایی گروه‌ها، دقت نمایش صحیح و تنظیم projection چطور نتیجه را عوض می‌کنند؟ کد sweep کنترل‌شده اجرا می‌کند تا loss داده‌ای، projection و quantization با هم اشتباه نشوند.

جدول ablation سؤال مکمل است: اگر یک جزء روش حذف یا جایگزین شود، کیفیت و هزینه چه تغییری می‌کند؟ عددهای symbolic برای proof جای اندازه‌گیری نیستند. variant پشتیبانی‌نشده در مدار، مثل \lr{incircuit-gen} فعلی، باید \lr{native-only baseline} باشد، نه ردیف proof با قید ساختگی.

\takeaway{RQ1 خوب، شرط لازم امید به روش است؛ برای ادعای برتری ZK هنوز RQ2 موفق و منصفانه لازم داریم.}

\lesson{RQ1: نتیجه فعلی، بهترین نقاط و محدودیت‌ها}
\status{Native؛ ۱۰۷۰ اندازه‌گیری، ۱۰۷ سلول پارامتری با ده seed. روی زیرمجموعه‌های جدول دیتاست‌ها.}
در مقایسه برابر با \(m=64,q=2\)، میانگین PCA در همه هفت خانواده از SparseJL بهتر بود. میله‌های شکل فاصله اطمینان ۹۵درصدیِ میانگین ده seed را نشان می‌دهند؛ این یک مقایسه تنظیمات مشخص است، نه حکم درباره همه داده‌های ممکن.
\chart[0.88]{quality_matched.png}{همه روش‌ها در این شکل بُعد خروجی یکسان دارند. مرجع، انتساب به مرکزهای فضای اصلی است. Full نسبت به خودش ARI برابر یک دارد.}

در grid بررسی‌شده، نقطه‌های کاهش‌یافته‌ای انتخاب کردیم که کران پایین CI برای ARI حداقل ۰٫۸ و تمام CI نسبت objective در بازه ۰٫۹ تا ۱٫۱ بود:
\begin{latin}\begin{center}\small
\begin{tabular}{lrrrr}
\toprule Dataset & \(d\to m\) & \(q\) & Mean ARI [95\% CI] & Reduction \\\midrule
Arcene & \(10000\to16\) & 2 & .944 [.907, .980] & 99.84\% \\
Fashion & \(784\to64\) & 2 & .848 [.833, .863] & 91.84\% \\
Gaussian & \(256\to32\) & 2 & .951 [.912, .990] & 87.50\% \\
MNIST & \(784\to256\) & 2 & .870 [.856, .883] & 67.35\% \\
\bottomrule\end{tabular}\end{center}\end{latin}
این‌ها نقاط انتخاب‌شده اکتشافی از همان داده‌اند، نه optimum از پیش ثبت‌شده یا تضمین روی داده جدید. برای Gisette، متن و synthetic DBSCAN هیچ نقطه واقعاً کاهش‌یافته با این شروط پیدا نشد. افزایش \(q\) از ۲ به ۸ هم بهبود ثابت و سراسری نداد.

اختلاف متوسط مطلق توافقِ مسیر اعشاری و صحیح حدود ۰٫۰۰۰۵۴۶ و بیشترین اختلاف یک ردیف حدود ۰٫۰۰۹۰۸۷ بود. این شواهد نشان می‌دهند در تنظیمات RQ1، افت اصلی غالباً از projection است، نه یک truncation شدید باقی‌مانده. ولی آزمایش native جای proof متصل به transcript را نمی‌گیرد.

\takeaway{روش بی‌فایده نیست؛ منطقه‌های مناسب دارد. اما در دقت بومی از PCA جلو نزده و روی متن و Gisette در بُعد کوچک ضعیف است.}

\lesson{RQ2: چرا این‌قدر سنگین است و PCA هم مشکل مشترک دارد؟}
\status{دو تلاش نماینده بزرگ قبلی ناموفق؛ proof کوچک آموزشی این فایل موفق. بنچمارک بزرگِ مقایسه‌ای تکمیل نشده است.}
سؤال تو درست است: وقتی PCA هم فشرده کرده، چرا هنوز داده اصلی سنگین را حساب می‌کنیم؟ چون باید ثابت شود خروجی از \textbf{همان داده متعهدشده} آمده است. فقط بررسی فاصله‌های فشرده این ارتباط را اثبات نمی‌کند. هر دو روش در تعریف فعلی باید داده اصلی را بخوانند و تعهد آن را احراز کنند.

سه هزینه را جدا ببین: \textbf{تعهد و ارتباط با داده اصلی} با \(nd\) رشد می‌کند؛ \textbf{اعمال تبدیل} به ساختار ماتریس وابسته است؛ \textbf{فاصله‌های کاهش‌یافته} حدود \(nkm\) مختصه دارند. هدف SparseJL عمدتاً کاهش بخش‌های تبدیل و فاصله است، اما اگر تعهد غالب باشد، سرعت کل محدود می‌ماند.

مثلاً برای \(n=2000,d=256,k=5,m=16\)، تعهد با ۵۱۲هزار مختصه ورودی سروکار دارد، در حالی که مقایسه فاصله‌ها به ۱۶۰هزار مختصه کاهش‌یافته مرتبط است. برای \(n=100,d=10000\) با همان \(k,m\)، یک میلیون مختصه ورودی و فقط ۸هزار مختصه فاصله داریم. کاهش \(m\) مشکل میلیون مختصه تعهد را حذف نمی‌کند.

در دو تلاش قبلی، watchdog حافظه نزدیک ۶٫۸ گیگابایت فرایند را متوقف کرد، قبل از کامل‌شدن compile. مقدار صفر در فیلد constraints ردیف شکست، اندازه‌گیری «مدار بدون قید» نیست؛ مقدار پیش‌فرض یک نتیجه ناموفق است.

برآورد ساختاریِ hashها بسیار بزرگ است، اما نه عدد اندازه‌گیری‌شده قیود است و نه اثبات این‌که حداقل ۲۵۶ گیگابایت لازم داریم. برخی backendها سقف دامنه FFT دارند؛ مدار بیش‌ازحد بزرگ ممکن است حتی با RAM بیشتر هم به بازطراحی نیاز داشته باشد. \textbf{زمان کامل و حداقل RAM همه gridها اکنون با قطعیت معلوم نیست.}

PCA عمومی می‌تواند ضرایب ثابت داشته باشد؛ ضرب در ثابت در R1CS با ضرب دو witness یک هزینه ندارد. SparseJL فعلی ماتریس شاهد و احراز دارد و routing آن هنوز بهینه نیست. بنابراین شمردن ضرب‌های اعشاری به‌تنهایی مقایسه معتبر اثباتی نیست.

\takeaway{نباید مزیت نهایی ZK را از کاهش مختصات یا RQ1 نتیجه بگیریم. باید overhead مشترک تعهد و تفاوت مدارهای واقعی را با proof موفق اندازه بگیریم.}

\lesson{راه‌های حل گلوگاه RQ2، بدون کوچک‌کردن ادعای مقاله}
اول باید profile کنیم: compile، hash/Merkle، routing تنک، range check، setup و prove هرکدام چه سهمی دارند؟ اگر bottleneck مشخص نشود، خرید RAM یا زیادکردن هسته ممکن است فقط هزینه را بالا ببرد.

\textbf{بهینه‌سازی ساختار مدار:} routing فعلی می‌تواند بسیار بیشتر از \(ndq\) عملیات منطقی بسازد. gadget مناسب برای sparse lookup/selection و بازبینی محدودیت‌ها می‌تواند کمک کند؛ باید با تست witness نامعتبر soundness حفظ شود. حذف constraint صرفاً برای سرعت، اگر رابطه را ضعیف کند راه‌حل نیست.

\textbf{تعهد کم‌هزینه‌تر یا معماری proof متفاوت:} سازوکارهای commitment-consistency، مانند مسیر پژوهشی CP-SNARK، برای کاهش همین overhead مطرح شده‌اند. Artemis یک نمونه پژوهشی مرتبط است [۱۰]، اما جایگزین آماده و پیاده‌سازی‌شده برای MiMC Merkle این مخزن نیست. GKR یا معماری‌های تجمیع هم گزینه تحقیق‌اند، نه قابلیت تمام‌شده فعلی.

\textbf{Batching دقیق:} می‌توان همه ردیف‌ها را در batchهای کوچک‌تر پردازش کرد، به شرط حفظ تعهد کلی، مرکزها، ماتریس، پوشش همه نمونه‌ها و جمع صحیح خروجی‌ها. باید از حذف، تکرار و جابه‌جایی batch جلوگیری شود. هزینه نهایی مجموع همه proofها و verificationهاست؛ زمان یک batch را نباید به‌عنوان full run گزارش کرد. این تغییر ممکن است رابطه و فرض‌های پروتکل را تغییر دهد و باید صریح نوشته شود.

\textbf{استفاده دوباره از setup:} در تکرارهای یک شکل سازگار، compile/setup را یک‌بار انجام می‌دهیم. شکل متفاوت معمولاً مدار متفاوت است. reuse مقدار اثبات‌های لازم را کم نمی‌کند؛ فقط تکرار کار مشترک را حذف می‌کند.

\textbf{سرور ۶۴ گیگابایتی:} برای nativeها بسیار مناسب‌تر از لپ‌تاپ است، اما کامل‌شدن همه proofهای بزرگ با آن تضمین‌شده نیست. باید شکل‌ها را تدریجی profile کنیم و مراقب سقف backend باشیم. بیان «تمام ران‌ها حتماً ۲۵۶ گیگابایت می‌خواهند» اکنون قابل دفاع نیست؛ آن عدد تخمین بود، نه حداقل تجربی.

\takeaway{راه علمی، ابتدا کوچک‌ترین proofهای نماینده موفق، سپس profile و بهینه‌سازی تعهد و مدار، و بعد گسترش با ادعای یکسان است. تغییر روش هنوز یک تصمیم باز است.}

\lesson{RQ3: سربار کم، ولی معنی امنیت باید دقیق باشد}
\status{۹۰ اجرای زمان‌سنجی با ده seed در ۹ سلول، و ۸ سلول شبیه‌سازی grinding. Native؛ beacon محلی شبیه‌سازی شده.}
برای \(d=128,m=32,q=2,k=3\) و \(n=64,128,256\)، کل پروتکل میزبان به‌طور متوسط حدود ۸٫۸۷، ۱۷٫۵۵ و ۳۶٫۵۵ ثانیه بود. مرحله پس از commit برای challenge و ساخت/احراز ماتریس حدود ۰٫۱۴۱ تا ۰٫۱۵۱ ثانیه شد. سهم آن از کل زمان حدود ۱٫۵۹، ۰٫۸۱ و ۰٫۴۱ درصد بود. سهم بزرگ مربوط به commit داده است، نه انتخاب منبع تصادف.

این timing proof ZK بزرگ را شامل نمی‌شود. اندازه‌ها هم اندازه کامل دیتاست‌های کیفیت نیستند. CI اختلاف paired کل تراکنش با و بدون اضافه مربوطه از صفر عبور می‌کند؛ پس جریمه کل زمان با این داده‌ها تفکیک آماری روشنی ندارد. این به معنی اثبات دقیقاً صفر بودن هزینه نیست.
\chart{overhead_clean.png}{میانگین توصیفی زمان برحسب اندازه؛ سمت راست سهم challenge و ساخت/احراز ماتریس از کل میزبان است. این اعداد هزینه شبکه beacon واقعی یا prove بزرگ را شامل نمی‌شوند.}

در grinding، Gaussian و MNIST با هزار نمونه بررسی شدند. برای هر سلول ۳۲ projection جهت تنظیم آستانه pilot و ۲۵۶ projection مستقل برای ارزیابی داریم. رویداد، عبور distortion جفت‌های نمونه‌گیری‌شده از آستانه است؛ نه پذیرفتن یک proof جعلی.

با آستانه مشترک \(\varepsilon=0.30,q=2,Q=1\)، در \(m=64\) نرخ عبور ۱۰۰ درصد بود؛ در \(m=256\) حدود ۲۳٫۴ درصد Gaussian و ۲۱٫۱ درصد MNIST شد. این اثر واقعی افزایش بُعد در رویداد مشترک است. در آستانه سخت ۰٫۱۵، همه سلول‌ها اشباع‌اند و مقایسه retry چندان آموزنده نیست.

در آستانه pilot جدا برای \(m=64,q=2\)، نرخ Gaussian از ۲۷٫۳ درصد در یک تلاش به ۸۷٫۵ درصد در هشت تلاش رسید؛ MNIST از ۲۷٫۷ به ۹۳٫۸ درصد. آستانه‌های pilot بین سلول‌ها فرق دارند؛ این درصدها را برای ادعای برتری یک بُعد بر دیگری مستقیم مقایسه نمی‌کنیم. \(Q=8\) فقط ۳۲ گروه مستقل دارد و CI آن پهن‌تر است.

\takeaway{تازگی transcript در محاسبه میزبان سربار کوچکی دارد و retry می‌تواند رویداد بد هندسی را بیشتر کند؛ هیچ‌کدام اثبات نرخ شکستن ZK نیستند.}

\lesson{RQ4: گواه درست، پوشش متفاوت و نقطه‌های سخت}
\status{۱۲۶۰ ردیف؛ ۱۲۰۰ ردیف گواه مرکز و ۶۰ ردیف همسایگی DBSCAN. Native، ده seed در هر سلول.}
روی مجموعه‌های eligible، تغییر مشاهده‌شده صفر بود. تعداد بررسی‌ها ۱۶۴۵۸۷ مشاهده نقطه‌ای و ۲۴۷۹۱۳۸۴ مشاهده جفتی بود؛ این‌ها روی پارامترها و seedهای مختلف تکرار می‌شوند. نباید آن‌ها را میلیون‌ها نمونه مستقل یا کران امنیتی برای شکست بسیار نادر حساب کنیم.

پوشش گواه مرکز در \(m=256,q=2,\varepsilon=0.15\): Arcene حدود ۷۱٫۴ درصد، Fashion حدود ۵۲٫۴، MNIST حدود ۳۲٫۱، متن حدود ۸ و Gisette حدود ۰٫۶۱ درصد. پس روی Gisette خیلی از نقطه‌ها اصلاً از این گواه کافی بهره نمی‌برند، حتی اگر بعضی انتساب‌ها اتفاقاً حفظ شوند.
\chart{coverage_clean.png}{میانگین پوشش و CI ساده t بین seedها؛ کران منفی CI نرخ منفی نیست. این میله‌ها صحت classifier یا تضمین تمام دیتاست نیستند.}

برای DBSCAN مصنوعی تمام ۴۹۹۵۰۰ جفت بررسی شد، ولی فقط ۲۵۰ جفت نزدیک مرز داریم. در همین جفت‌های سخت، پایداری متوسط از ۵۸٫۴ درصد با CI حدود [۵۶٫۲،۶۰٫۷] در \(m=32\) به ۶۵٫۸ درصد با CI [۶۳٫۸،۶۷٫۷] در \(m=128\) افزایش یافت.

این نمونه خوبی از دام میانگین کلی است: چون اکثر جفت‌ها خیلی دورند، ممکن است درصد کلی بسیار بالا باشد، ولی درصد قابل توجهی از تصمیم‌های نزدیک شعاع جابه‌جا شوند. برای کاربرد وابسته به اتصال خوشه‌ها، همین بخش کوچک می‌تواند اثر مهمی داشته باشد.

صفر تغییر در eligible یعنی پیاده‌سازی شرط کافی در این بررسی با مشاهده سازگار است. برای تضمین عمومی همچنان باید رویداد هندسی، همه بردارهای لازم و قرارداد الگوریتم تعریف شوند. ردیف‌هایی با eligible صفر هم «۱۰۰ درصد درست» نیستند؛ نرخ آن مجموعه تعریف‌نشده است و باید جدا نمایش داده شود.

\takeaway{نتیجه RQ4 از استدلال شرطی حمایت می‌کند، اما ثابت نمی‌کند همه تصمیم‌ها یا کل خوشه‌بندی DBSCAN در تمام ران‌ها حفظ شده‌اند.}

\lesson{RQ5: جدایی خوشه‌ها، دقت صحیح و یک کنترل مهم}
\status{۸۵۰ ردیف: ۷۰۰ برای دقت روی هفت دیتاست و دو سیاست مقیاس، ۱۵۰ برای جدایی روی داده مصنوعی.}
وقتی خوشه‌ها جدا باشند، projection کوچک هم تصمیم را بهتر نگه می‌دارد. برای \(m=64,q=2\)، توافق با مرجع در \(\gamma=0.1,0.25,0.5,1,2\) تقریباً ۰٫۱۴۱، ۰٫۴۶۵، ۰٫۹۲۲، ۰٫۹۹۹ و ۱ شد.
\chart{separation_clean.png}{اثر جدایی با سه بُعد خروجی و CI میانگین ده seed. این نتیجه از مولد فعلی است، نه تعریف تمام مسئله‌های واقعی.}
در \(\gamma=0.1\)، حتی مرجع اصلی نسبت به برچسب مولد ARI حدود ۰٫۰۱۹ دارد؛ خود داده سخت است. در \(\gamma=0.25\)، مرجع اصلی نسبت به label حدود ۰٫۸۹۸ دارد ولی کاهش بُعد افت قابل توجهی می‌آورد؛ این‌جا loss روش روشن‌تر دیده می‌شود.

در sweep دقت، برچسب‌های ۸، ۱۲، ۱۶، ۲۰ و ۲۴ بیت و \(m=32,q=2\) داریم. سیاست اول یک مرجع فیزیکی ۲۴بیتی ثابت نگه می‌دارد. در آن، ۸بیت همه مختصات را صفر می‌کند؛ متن در ۱۶بیت ARI حدود ۰٫۰۳۷ و در ۲۴بیت حدود ۰٫۰۹۴ دارد. MNIST حدود ۰٫۵۴۱ و ۰٫۵۴۴؛ Arcene حدود ۰٫۹۶۷ و ۰٫۹۷۴ دارد.
\chart{precision_clean.png}{میانگین ARI در سیاست مرجع ثابت؛ این شکل را نباید به همه شیوه‌های مقیاس‌دهی تعمیم داد.}

سیاست کنترل، دامنه مجاز \(B\) را با مقیاس‌دهی دوباره پر می‌کند. در آن، برای هر دیتاست و seed، همه پنج برچسب دقت همان توافق صحیح را می‌دهند؛ عدد صحیح می‌تواند یکسان باشد. بنابراین گزاره «۸بیت همیشه شکست می‌خورد» غلط است. همچنین کران مدار ثابت مانده، پس «۱۶بیت optimum هزینه قید است» هنوز ثابت نشده است.

\takeaway{دقت را باید با witness واقعی و range check تعریف کنیم، نه صرفاً اسم bits. loss مقیاس، loss projection و سختی خود داده را جدا می‌کنیم.}

\lesson{آمار ساده: ده seed و فاصله اطمینان چه می‌گویند؟}
یک بار اجرا ممکن است خوش‌شانس یا بدشانس باشد. ده seed برای دیدن تغییرات بین اجراها استفاده می‌شود. برای معیار \(y\)، میانگین \(\bar y\) مرکز نتایج و انحراف معیار پراکندگی خود اجراها را توصیف می‌کند. فاصله اطمینان میانگین با ده مقدار تقریباً:
\[
 \bar y\ \pm\ t_{0.975,9}\frac{s}{\sqrt{10}},\qquad t_{0.975,9}\approx2.262.
\]
این CI درباره میانگین تحت مفروضات آماری است، نه بازه‌ای که حتماً ۹۵ درصد اجرای آینده داخلش می‌افتد. CI همچنین نااطمینانی انتخاب دیتاست یا یک نمونه دیگر از کل جمعیت را خودکار شامل نمی‌شود.

در برخی سلول‌ها همان seed هم initialization KMeans و هم projection را تغییر می‌دهد. این طراحی کل تغییرات مشترک را می‌سنجد، نه اثر مستقل یکی را به‌تنهایی. نمونه‌گیری دیتاست در بسیاری از sweepها ثابت است. بنابراین ده seed یعنی ده تکرار روی همان instance، نه ده دیتاست مستقل.

برای مقایسه paired، خروجی دو تنظیم روی همان seed را کم می‌کنیم و CI اختلاف را می‌سازیم؛ معمولاً واضح‌تر از مقایسه دو CI جداست. اگر CI اختلاف از صفر عبور کند، با این داده‌ها تفاوت قطعی گزارش نمی‌کنیم؛ نبود شواهد تفاوت، اثبات یکسان‌بودن نیست.

وقتی ده‌ها تنظیم می‌سنجیم و بهترین را انتخاب می‌کنیم، انتخاب پس از مشاهده می‌تواند خوش‌بینانه باشد. نقاط منتخب RQ1 اکتشافی‌اند. در تفاوت‌های precision فعلی، CIها بدون تصحیح چندمقایسه‌ای گزارش شده‌اند؛ آن‌ها را significance نهایی و تأییدشده مقاله نمی‌نامیم.

برای نرخ grinding از گروه‌های مستقل و CI مناسب نرخ استفاده می‌شود. صفر رویداد از ۲۵۶ draw یعنی «در این ۲۵۶ مورد ندیدیم»، نه احتمال واقعاً صفر؛ کران بالای Wilson حدود ۱٫۵ درصد می‌تواند باقی بماند.

\takeaway{ده seed کمک می‌کند نتیجه تصادفی یک اجرا را نادیده نگیریم، ولی جای طراحی منصفانه، پوشش داده و اعتبارسنجی مستقل را نمی‌گیرد.}

\lesson{روش ما کجا ارزش دارد، کجا ضعیف است و چرا PCA را کنار نمی‌گذاریم؟}
\textbf{جایی که شواهد فعلی امیدوارکننده‌اند:} داده‌های با جدایی مناسب، Arcene برای حفظ مرجع همین نمونه، Fashion در بُعد ۶۴، و MNIST با بُعد بزرگ‌تر ۲۵۶. ساختار علامت‌دار تنک و seed پس از تعهد، برای یک پروتکل قابل احراز جذاب‌اند. مرحله تازگی روی میزبان سهم زمانی کوچکی دارد.

\textbf{جایی که شواهد ضعیف‌اند:} کاهش شدید بُعد روی متن و Gisette، جفت‌های DBSCAN نزدیک شعاع، و خوشه‌های مصنوعی کم‌جدایی. ARI بالا در Arcene کاربرد پزشکی را ثابت نمی‌کند. حفظ objective نزدیک یک هم حفظ انتساب را تضمین نمی‌کند. برای \(\varepsilon=0.15\) رویداد خطای سخت در بعضی تنظیم‌ها برقرار نیست.

\textbf{چرا کسی ممکن است SparseJL را به PCA ترجیح دهد؟} ماتریس مستقل از یادگیری جهت‌های داده است، تعداد ناصفر مشخص دارد و ضرایب صحیح علامت‌دارند؛ می‌توان قواعد تولید پس از commit را تعریف کرد. این مزایا برای امنیت ترتیب، نمایش داده و معماری اثبات \textbf{فرضیه کاربردی} می‌سازند. اما انتخاب بر پایه سرعت ZK تا تکمیل RQ2 موفق، اثبات تجربی ندارد.

\textbf{حوزه‌هایی که باید هدفمند آزمایش شوند:} تحلیل خصوصی با مرکزهای عمومیِ از قبل تعیین‌شده؛ دیتاست‌های کم‌نمونه و پُرویژگی با ساختار قابل کاهش؛ مسئله‌های با margin مناسب؛ و سناریوهای چند پرس‌وجوی سازگار که ممکن است هزینه مشترک احراز داده را توزیع کنند. مورد آخر در نسخه فعلی به‌عنوان کاهش هزینه کامل اثبات نشده است.

کاربردهای مالی، پزشکی یا embeddingهای زبانی، انگیزه‌اند نه نتیجه فعلی. برای هر کاربرد باید معیار واقعی خودش هم سنجیده شود: تصمیم پایین‌دستی، label واقعی، هزینه اشتباه و حریم خصوصی خروجی. روش صرفاً چون «ZK-friendly» نامیده شود برتری کاربردی ندارد.

\takeaway{انتخاب منطقی این است: اگر fidelity بومی اولویت است، PCA فعلاً مرجع قوی‌تر است. اگر پروتکل خصوصی و قابل احراز هدف است، SparseJL ارزش بررسی دارد، با RQ2 هنوز باز.}

\lesson{مقاله از ابتدا تا انتها چه روایت می‌کند؟}
\textbf{مقدمه:} مشکل محاسبه خصوصی و سنگینی proof را طرح می‌کند؛ کاهش بُعد را راه کاهش بخشی از هزینه معرفی می‌کند و خطر انتخاب داده پس از ماتریس را نشان می‌دهد. سهم پیشنهادی، اتصال projection تنک به transcript پس از تعهد و نتیجه‌های پایداری است.

\textbf{کارهای مرتبط:} روش‌های projection و JL، سامانه‌های ZK برای محاسبات و مسیرهای commitment-consistency را کنار هم می‌گذارد. مقاله ezDPS مثلاً از DWT و PCA در pipeline یادگیری خصوصی استفاده می‌کند [۹]؛ مسئله آن دقیقاً خوشه‌بندی ثابت‌مرکز ما نیست. مقاله Pisot موضوع projection قطعی برای private routing دارد [۱۱]؛ baseline فعلی ما بازتولید مستقل تأییدشده artifact آن نیست.

\textbf{تعریف‌ها و مدل تهدید:} داده، مرکزها، ماتریس، پارامترها و آنچه عمومی یا خصوصی است را تعریف می‌کند. ترتیب انتخاب، بودجه retry و نقش منبع تصادف باید واضح باشد. هر فرضی که فقط روی میزبان enforce می‌شود، باید به verifier کاربردی نسبت داده شود، نه به خود مدار.

\textbf{پروتکل و رابطه مدار:} زنجیره commit تا challenge و اثبات صحیح‌بودن محاسبه را توصیف می‌کند. کد فعلی مسیر field-native دارد؛ بعضی جمله‌های پیش‌نویس درباره mismatch encoding قدیمی‌اند. پروتکل روی همان نسخه پیاده‌سازی‌شده باید بازنویسی شود، نه این‌که خطای قدیمی به‌عنوان وضعیت امروز باقی بماند.

\textbf{تحلیل نظری:} از رویداد حفظ فاصله به کران هزینه، margin انتساب و مرز همسایگی می‌رسد. این‌ها نتیجه شرطی‌اند؛ پارامترهای عملی باید جدا نشان دهند شرایط کجا برقرار یا برقرار نیستند.

\textbf{آزمایش‌ها و بحث:} پنج RQ می‌پرسند کیفیت، هزینه، امنیت ترتیب، پوشش گواه و حساسیت چگونه‌اند. پیش‌نویس هنوز جاهای نتیجه/طرح باز دارد. نتایج تشخیصی این فایل می‌توانند جهت بحث را روشن کنند، اما همه جدول‌های کامل مقاله را پر نمی‌کنند. بحث باید محدودیت‌های native-only، زیرمجموعه‌ها و RQ2 ناموفق را صریح حفظ کند.

\takeaway{روایت قابل دفاع «روش همیشه بهتر است» نیست؛ «یک طراحی قابل احراز داریم، محدوده کیفیت آن را شناختیم و گلوگاه اثبات آن هنوز باید حل و اندازه‌گیری شود» است.}

\lesson{نقشه کد و بازتولید مثال‌ها}
کد اصلی الگوریتم و پروتکل در package محلی \lr{commitjl} است. بارگذارها در \lr{data}، gridها و pipeline آزمایش در \lr{experiments}، اجراکننده کمپین‌ها در \lr{campaigns} و کد تحلیل در \lr{analysis} قرار دارند. مدار و CLI اثبات زیر بخش \lr{zk} هستند. نام دقیق entry pointها را از پوشه‌ها و راهنمای جاری مخزن بخوانیم؛ تغییر ساختار اخیر بعضی نام‌های قدیمی را عوض کرده است.

برای این فایل، یک اسکریپت کوچک آموزشی نوشته و اجرا شده است. هیچ نتیجه اصلی بنچمارک بازنویسی نشده و یک full campaign خودکار شروع نشده است. خروجی‌ها از نتایج مقاله جدا و در پوشه آموزش ذخیره شده‌اند.
\begin{latin}\begin{lstlisting}[language={}]
# Run from PowerShell in the project root:
Set-Location "C:\PyProjects\zero-knowledge"
python code/analysis/persian_guide_examples.py
python code/analysis/build_persian_guide.py
\end{lstlisting}\end{latin}
دستور اول به داده‌های محلی و dependencyهای علمی و bridge Go نیاز دارد؛ نمونه کوچک proof هم می‌سازد. دستور دوم برای بازسازی PDF به XeLaTeX/XePersian، فونت Tahoma و کتابخانه‌های plotting/PDF نیاز دارد. این‌ها دستورهای full paper benchmark نیستند.

\begin{latin}\begin{lstlisting}[language={}]
manuscript.tex
code/analysis/persian_guide_examples.py
code/analysis/commitjl_guide_fa.tex
code/analysis/build_persian_guide.py
code/results/overview/education/persian_guide/examples.json
code/results/overview/education/persian_guide/document_manifest.json
code/results/diagnostics/local_v1/measurements.parquet
code/results/diagnostics/local_v1/analysis/reviewed_findings.md
output/pdf/commitjl_complete_guide_fa.pdf
\end{lstlisting}\end{latin}
JSON مثال‌ها شامل نمونه ویژگی‌ها، شکل داده، خروجی PCA، transcript و اطلاعات proof کوچک است. manifest سند، hash منبع و اندازه‌گیری‌ها را ثبت می‌کند تا بدانیم گزارش از کدام نسخه تولید شده است. برای نگاه شهودی، شکل دوبعدی PCA صرفاً نمایش است؛ آزمایش خوشه‌بندی واقعی در دو بُعد آموزش نشده است.
\chart{mnist_pca_plane.png}{نمایش دو مؤلفه اول؛ رنگ‌ها پنج خوشه فضای اصلی هستند، نه ده کلاس رقم.}

\lesson{سؤال‌های جلسه: جواب کوتاه و دقیق}
\textbf{آیا روش ما دقت دارد؟} روی بعضی تنظیمات، انتساب مرجع را خوب حفظ می‌کند؛ روی بعضی‌ها نه. PCA در مقایسه برابرِ بُعد ۶۴ بهتر بود. «دقت» باید بگوید نسبت به مرجع خوشه‌ای است یا label واقعی.

\textbf{آیا ثابت کردیم در ZK سبک‌تر است؟} نه، proofهای بزرگ نماینده ناموفق بوده‌اند. نمونه کوچک واقعی موفق است، اما مزیت مقایسه‌ای بزرگ هنوز اندازه‌گیری نشده.

\textbf{آیا native یعنی داده کوچک؟} نه؛ native یعنی بدون proof. ران فعلی روی اندازه‌های مشخص زیرمجموعه/instance انجام شده است. \textbf{آیا full را گرفته‌ایم؟} تمام ردیف‌های کمپین تشخیصی تمام شده‌اند، اما تمام طراحی بزرگ مقاله تمام نشده است.

\textbf{آیا کم‌کردن \(m\) تمام هزینه را کم می‌کند؟} خیر؛ احراز داده اصلی با \(nd\) باقی می‌ماند و routing/تعهد می‌توانند غالب باشند. \textbf{آیا RAM آزاد یعنی برنامه ناکارآمد است؟} نه؛ بعضی مرحله‌ها تک‌فرایندی، dependencyدار یا محدود به حافظه‌پهنای‌باندند. پرکردن RAM هدف نیست؛ کم‌کردن زمان با حفظ اندازه‌گیری معتبر هدف است.

\textbf{آیا ۶۴ گیگابایت کافی است؟} برای nativeها مسیر عملی قوی‌تری است؛ برای همه proofهای کامل هنوز تضمینی نداریم. \textbf{آیا حتماً ۲۵۶ لازم است؟} حداقل تجربیِ ثابت‌شده نداریم؛ عدد قبلی تخمین بوده. سقف backend هم باید بررسی شود.

\textbf{آیا ۸بیت خراب و ۱۶ بهترین است؟} نه به‌طور عمومی. سیاست مقیاس مهم است و کنترل range-utilized آن برداشت ساده را رد می‌کند. \textbf{آیا گواه ۱۰۰درصد یعنی همه نقاط امن‌اند؟} نه؛ باید پوشش eligible را کنار correctness آن بخوانیم.

\textbf{آیا مقاله را همین حالا باید کنار گذاشت؟} داده‌ها دلیل کنارگذاشتن کامل ایده نیستند، ولی ادعای گسترده برتری هم پشتیبانی نمی‌شود. تصمیم بهتر این است که دامنه کاربرد را محدود و شفاف کنیم و RQ2 را با معماری و baseline منصفانه به نتیجه واقعی برسانیم.

\takeaway{در جلسه، قوت واقعی روش را بیان کن، ولی فاصله بین فرضیه، پیاده‌سازی و اندازه‌گیری را هیچ‌جا پنهان نکن.}

\lesson{واژه‌نامه و منابع برای ادامه مطالعه}
\textbf{Feature:} یک ستون داده. \textbf{Embedding/Projection:} نمایش عددی کوتاه‌تر؛ در این پروژه عمدتاً تبدیل خطی. \textbf{Centroid:} مرکز یا میانگین یک خوشه. \textbf{Margin:} فاصله تصمیم اول از رقیب نزدیک. \textbf{Sparse:} دارای صفرهای زیاد. \textbf{Dense:} عمدتاً ناصفر. \textbf{Witness:} داده خصوصی لازم برای ارضای رابطه اثبات. \textbf{Statement:} ادعای عمومی دقیق. \textbf{Constraint:} یک رابطه که witness باید ارضا کند. \textbf{Commitment:} تعهد به مقدار. \textbf{Transcript:} زنجیره پیام‌ها و مقادیر پروتکل. \textbf{Seed:} ورودی تکرارپذیر مولد. \textbf{Soundness:} مقاومت در برابر پذیرفتن ادعای غلط. \textbf{Grinding:} تکرار انتخاب برای رسیدن به نتیجه مطلوب. \textbf{Quantization:} نمایش عدد واقعی با شبکه صحیح محدود. \textbf{CI:} فاصله اطمینان آماری. \textbf{Native-only:} اندازه‌گیری بدون ساخت proof مدار.

منابع اصلی آموزشی و پژوهشی زیر برای تعریف‌ها و توصیف دیتاست‌ها استفاده شدند؛ اعداد نتایج این فایل از خروجی محلی خود پروژه آمده‌اند، نه از مقاله‌های دیگر. توضیح‌ها و مثال‌های محاسباتی بازنویسی آموزشی‌اند.
\begin{latin}\small\linespread{1.05}\selectfont
\begin{enumerate}
\item Scikit-learn, PCA documentation.\par
\url{https://scikit-learn.org/stable/modules/generated/sklearn.decomposition.PCA.html}
\item Scikit-learn, KMeans documentation.\par
\url{https://scikit-learn.org/stable/modules/generated/sklearn.cluster.KMeans.html}
\item Consensys gnark, Zero-knowledge proofs concepts.\par
\url{https://docs.gnark.consensys.io/Concepts/zkp}
\item UCI Machine Learning Repository, Gisette.\par
\url{https://archive.ics.uci.edu/dataset/170/gisette}
\item Zalando Research, Fashion-MNIST original repository.\par
\url{https://github.com/zalandoresearch/fashion-mnist}
\item UCI Machine Learning Repository, Arcene.\par
\url{https://archive.ics.uci.edu/dataset/167/arcene}
\item Scikit-learn, 20 Newsgroups dataset guide.\par
\url{https://scikit-learn.org/stable/datasets/real_world.html}
\item Kane and Nelson, Sparse Johnson--Lindenstrauss Transforms.\par
\url{https://arxiv.org/abs/1012.1577}
\item Wang and Hoang, ezDPS: An Efficient and Zero-Knowledge Machine Learning Inference Pipeline (2023).\par
\url{https://petsymposium.org/popets/2023/popets-2023-0061.php}
\item Artemis: Efficient Commit-and-Prove SNARKs for zkML.\par
\url{https://arxiv.org/abs/2409.12055}
\item Dey and Cherkaoui, Deterministic Johnson--Lindenstrauss Projections from Pisot beta-Transformations for Zero-Knowledge Private Routing (2026).\par
\url{https://arxiv.org/abs/2608.13078}
\item Scikit-learn, DBSCAN documentation.\par
\url{https://scikit-learn.org/stable/modules/generated/sklearn.cluster.DBSCAN.html}
\end{enumerate}
\end{latin}
\textbf{منابع داخلی:} پیش‌نویس \lr{manuscript.tex}، کد فعلی پروتکل و مدار، نتایج \lr{local\_v1}، یادداشت یافته‌های بازبینی‌شده و \lr{examples.json}. نمونه‌های آموزشی عمداً جدا از نتیجه‌های بنچمارک‌اند. تغییرات احتمالی بعدی کد یا داده می‌توانند نیازمند بازسازی این سند باشند.
\end{document}
